\section{Introduction}

Moduli spaces of sheaves of fixed topological type that are Gieseker-Maruyama-semistable with respect to a given ample class on a projective manifold $X$ have been studied for several decades. When one studies the way these moduli spaces vary if the polarisation changes, examples show that in dimension bigger than two one encounters sheaves $\mathcal{E}$ that are Gieseker-Maruyama-semistable with respect to non-rational, real ample classes $\alpha \in \mathrm{Amp}(X)_\mathbb{R}$ on $X$, \ie that enjoy the property that for a Kähler form $\omega$ representing $\alpha$ and for every proper coherent subsheaf $\mathcal{F} \subset \mathcal{E}$ we have $p_{\mathcal{F}}(m) \leq p_\mathcal{E}(m)$ for all $m$ sufficiently large, where the reduced Hilbert polynomial $p_E(m)$ with respect to $\alpha = [\omega]$ is defined by
\[p_E(m) = \frac{1}{\rank(E)}\int_X \ch(E) e^{m\omega} \Todd(X),\]see for example \cite{GRTsurvey}. When $\omega$ represents the first Chern class of an ample line bundle $L$, the Riemann--Roch theorem states that $p_{E}(m)$ equals $\frac{1}{\rank(E)} \chi(E\otimes L^m)$, and so the above generalises the notion of Gieseker-Maruyama--stability from integral classes to real ample classes, and in fact to all Kähler classes $[\omega]$. Both in the case of a real ample polarisation and of an arbitrary Kähler class on a compact Kähler manifold, the question arises whether there is a moduli space for such sheaves. In fact, it seems that the problem of constructing such moduli spaces was explicitly posed quite some time ago by Tyurin, see the discussion in \cite[\S 3.2]{TelemanCommContMath}.

When semistability is measured with respect to an ample line bundle, the construction of moduli spaces is based on Geometric Invariant Theory and is hence of global nature. Using the special structure of cones of positive classes and Geometric Invariant Theory for moduli spaces of quiver representations, it was shown by the authors in joint work \cite{GRT1} with Julius Ross that a GIT-construction of projective moduli spaces for $\omega$-semistable sheaves can still be carried out on projective threefolds. When dealing with arbitrary compact Kähler manifolds it is however quite unlikely that a finite-dimensional, global construction of a moduli space is possible. As an alternative approach, it is natural to study the symmetries induced by automorphism groups on semi-universal deformation spaces and to carry out a functorial local construction from which in the end the moduli space is glued. This approach is most naturally pursued in the language of analytic/algebraic stacks. Using recent advances in this theory, both regarding the correct type of moduli space to construct \cite{Alper13} and regarding existence criteria \cite{AlperFS}, in this paper we establish the following main result:

\begin{thm*} The algebraic stack of $\omega$-semistable sheaves of rank two and given Chern classes admits a good moduli space that is a proper algebraic space; in particular, the moduli space is separated.
\end{thm*}

We do not expect the restriction to the rank two case to be necessary; here, it simplifies the analysis of the local slice models describing the action of the automorphism groups of stable sheaves on their semi-universal deformation space. Note however that the theorem stated above does not claim that the moduli space is projective or even a scheme; new methods seem to be needed to investigate these additional questions.

While the approach followed here is very promising in the general Kähler case, both fundamental work extending \cite{AlperFS} to the analytic setup and a finer analysis of the geometry of the symmetries of semi-universal analytic deformation spaces will be needed to attack the existence question for semistable sheaves on compact Kähler manifolds.

\subsubsection*{Structure of the paper}
In Section~\ref{sect:prelim}, we collect the basic notions and their fundamental properties. More precisely, Section~\ref{subsect:sheafextensions} discusses sheaf extensions and their automorphisms, in Section~\ref{subsect:coherentsheaves} we introduce the notion of Gieseker-Maruyama-semistability with respect to a Kähler class and establish the basic properties of this notion, in Section~\ref{subsect:rank2} we provide the structure theory of semistable sheaves of rank two, and in Section~\ref{subsect:quotient_presentations} we establish the fundamental geometric properties of the stack of semistable sheaves, with particular emphasis on local quotient presentations and slice models. In Section~\ref{sect:construction} the existence of a good moduli space is established by checking the conditions given in \cite[Th.\,1.2]{AlperFS}. In the final section, Section~\ref{sect:properties}, we identify the points of the moduli space as representing $S$-equivalence classes of sheaves and establish separatedness and properness of the moduli space, completing our investigation.

\subsubsection*{Acknowledgements}
The authors would like to thank Jarod Alper, Jochen Heinloth, and David Rydh for discussions regarding algebraic stacks and good moduli spaces. Moreover, they want to express their deep gratitude to Peter Heinzner who in the early stages of the project invited MT to Bochum twice and contributed to the discussions that lead to the development of the approach pursued here. Finally, they want to thank the referee for finding a gap in an earlier version of the proof of Corollary~\ref{cor:separation} and for a number of questions that helped to improve the overall readability of the paper.

\section{Basic notions and first properties}\label{sect:prelim}

\subsection{Global conventions}
We work over the field of complex numbers. All manifolds are assumed to be connected. We will work on a fixed complex projective manifold $X$ endowed with a cohomology class $\alpha \in H^{1,1}(X, \mathbb{R})$ that can be represented by a Kähler form $\omega$; \ie $\alpha = [\omega]$.

\subsection{Sheaf extensions and automorphisms}\label{subsect:sheafextensions}
Here we recall a few facts about sheaf extensions and state two lemmata to be used later in the paper.
We start by considering extensions of $\cO_X$-modules over a ringed space $(X,\cO_X)$, where $\cO_X$ is a sheaf of $\C$-algebras. It is known that the $\C$-vector space $\E(E_2,E_1)$ of classes of extensions of $E_2$ by $E_1$ modulo Yoneda equivalence is canonically isomorphic to $\Ext^1_{\cO_X}(E_2,E_1)$, \cf \cite[Exer.\,III.6.1]{Hartshorne}, \cite[Exer.\,A3.26]{Eisenbud_CommAlg}.
Morphisms $\alpha\in \Hom_{\cO_X}(E_1,E'_1)$, $\beta\in\Hom_{\cO_X}(E'_2,E_2)$ induce natural linear maps
$\alpha_*:\E(E_2,E_1)\to\E(E_2,E'_1)$,
$\beta^*:\E(E_2,E_1)\to\E(E'_2,E_1)$, \cf \cite[Exer.\,A3.26]{Eisenbud_CommAlg}. On the $\Ext^1$-side these correspond exactly to the linear maps induced by $\alpha$ and $\beta$ using the natural morphisms $\alpha_*:\Hom_{\cO_X}(E_2,E_1)\to\Hom_{\cO_X}(E_2,E'_1)$, $\beta^*:\Hom_{\cO_X}(E_2,E_1)\to\Hom_{\cO_X}(E'_2,E_1)$. It follows that $\alpha_*\circ\beta^*=\beta^*\circ\alpha_*$ in $\Hom_\C(\E(E_2,E_1),\E(E'_2,E'_1))$.

\begin{rem}\label{remark:extensions}
The following particular cases of the above construction will be used in the sequel:
\begin{enumerate}
\item\label{remark:extensions1}
When $E'_1=E_1$, $E'_2=E_2$, we get a natural action of
$\Aut(E_1)\times\Aut(E_2)$ on $\E(E_2,E_1)$ by\vspace*{-4pt}
\[
(\alpha,\beta)(\xi):= (\alpha_*\circ(\beta^{-1})^*)(\xi)=((\beta^{-1})^*\circ\alpha_*)(\xi),
\]
for $\alpha\in \Aut(E_1)$, $\beta\in\Aut(E_2)$, $\xi\in\E(E_2,E_1)$, \cf \cite[Chap.\,7]{LePotier_Lectures}. If moreover $E_1=E_2=:E$, we get a natural action of $\Aut(E)$ on $\E(E,E)$ by\vspace*{-2pt}
\[
\alpha(\xi):=(\alpha,\alpha)(\xi)= (\alpha_*\circ(\alpha^{-1})^*)(\xi)=((\alpha^{-1})^*\circ\alpha_*)(\xi).
\]
\item\label{remark:extensions2}
By functoriality, when $j\in \Hom_{\cO_X}(E_1,E'_1)$ admits a retract or when
$p\in\Hom_{\cO_X}(E'_2,E_2)$ admits a section, we get injective maps $j_*:\E(E_2,E_1)\to\E(E_2,E'_1)$,
$p^*:\E(E_2,E_1)\to\E(E'_2,E_1)$.
\end{enumerate}
\end{rem}

We next show how these considerations apply to infinitesimal deformations of sheaves. For simplicity we restrict ourselves to the case when $X$ is a compact analytic space and the sheaves involved are coherent, but note that similar arguments work in the category of coherent sheaves over schemes. We denote by $\nearrow:=(\mathrm{point}, \C[t])$ the double point, where $\C[t]:=\C[T]/(T^2)$ is the algebra of dual numbers over $\C$. Let~$F$ be a coherent sheaf on $X$ and $(S,0)$ a germ of a complex space. A deformation of $F$ with base $S$ is a pair $(\cF,\phi)$ where $\cF$ is a coherent sheaf on $X\times S$ flat over $S$ and $\phi:\cF_0\to F$ is an isomorphism. Two deformations $(\cF,\phi)$, $(\cF',\phi')$ of $F$ with base $S$ are called isomorphic if there exists an isomorphism of sheaves $\Phi:\cF\to\cF'$ such that $\phi'\circ\Phi=\phi$,
\cite[\S 4.2.2]{Palamodov-versal}.
There is a natural bijection between the set of isomorphism classes of deformations of $F$ with base $\nearrow$ also called (first-order deformations) and the vector space $\E(F,F)$, \cite[Th.\,2.7]{Hartshorne_Deformations}.
Any deformation of $F$ with base $S$ gives rise to a ``tangent map'' $T_0S\to\E(F,F)$. Finally we mention that the automorphism group of $F$ naturally acts on the set of (isomorphism classes of) deformations of $F$ with base $S$ by $g(\cF,\phi):=(\cF,g\circ\phi)$, for $g\in\Aut(F)$.

\begin{lemma}\label{lemma: deformations}
The natural identification between the set of isomorphism classes of first-order deformations of $F$ and $\E(F,F)$ is $\Aut(F)$-equivariant.
\end{lemma}

Fix now two coherent sheaves $E_1$, $E_2$ on $X$. In our set-up $W:=\E(E_2,E_1)$ is a finite dimensional complex vector space and there exists a universal extension\vspace*{-3pt}
\begin{equation}
0\to E_{1,W}\to \cE\to E_{2,W}\to 0
\end{equation}
on $X\times W$,
\cite[Chap.\,7]{LePotier_Lectures}. The central fibre of the universal extension is a trivial extension\vspace*{-5pt}
\begin{equation}
0\to E_{1}\To{\alpha} \cE_0\To{\beta} E_{2}\to 0
\end{equation}
on $X$.
Fixing a section $s\in \Hom_{\cO_X}(E_2,\cE_0)$ gives us an isomorphism $\phi:\cE_0\to E_1\oplus E_2$ hence a deformation $(\cE,\phi)$ of $E_1\oplus E_2$ with base $(W,0)$.

\begin{lemma}\label{lemma:extensions-deformations}
The tangent map $\E(E_2,E_1)\to \E( E_1\oplus E_2, E_1\oplus E_2)$ to the deformation $(\cE,\phi)$ induced by the universal extension coincides with the natural inclusion $(\phi\circ \alpha)_*\circ (\beta\circ\phi^{-1})^*=(\beta\circ\phi^{-1})^*\circ(\phi\circ \alpha)_*$ given by Remark \ref{remark:extensions}\eqref{remark:extensions2} and is equivariant with respect to the group homomorphism $\Aut(E_1)\times\Aut(E_2)\to\Aut(E_1\oplus E_2)$ and the actions described in Remark \ref{remark:extensions}\eqref{remark:extensions1}.
\end{lemma}
\begin{proof}
We will check that the images in $\E( E_1\oplus E_2, E_1\oplus E_2)$ of the class $\xi\in\E(E_2,E_1)$ of any extension
\begin{equation}
\label{equation:xi}
0\to E_{1}\To{j} E\To{p} E_{2}\to 0
\end{equation}
of coherent sheaves on $X$ induced in the two different ways described in the statement coincide. The second part of the lemma will follow from this.

Consider in addition a trivial extension
\begin{equation}
\label{equation: trivial extension}
0\to E_{1}\To{\alpha} E_0\To{\beta} E_{2}\to 0
\end{equation}
and fix a section $s:E_2\to E_0$ of $\beta$ and the induced retraction $r:E_0\to E_1$ of $\alpha$. Then it is directly seen that the class $\alpha_*(\xi)\in \E(E_2,E_0)$ is represented by the second line of the following commutative diagram:
\begin{equation*}
\begin{gathered}
\xymatrix@C=1.2cm{
0\ar[r] & E_1 \ar[d]_{\alpha}\ar[r]^-{j}
& E\ar[r]^-p \ar[d]^{\begin{smallpmatrix} \id\\ 0
\end{smallpmatrix}} & E_2 \ar[d]^{\id} \ar[r] & 0 \\
0\ar[r] & E_0 \ar[r]_-{\begin{smallpmatrix} j\circ r\\ \beta
\end{smallpmatrix}} & E\oplus E_2 \ar[r]_-{\begin{smallpmatrix} p&0
\end{smallpmatrix}} & E_2 \ar[r] & 0
}
\end{gathered}
\end{equation*}
and that the first line of the diagram
\begin{equation*}
\begin{gathered}
\xymatrix@C=1.5cm{
0\ar[r]
& E_0 \ar[d]_{\id}\ar[r]^-{\begin{smallpmatrix} j\circ r\\ s\circ\beta\end{smallpmatrix}}
& E\oplus E_0\ar[r]^-{\begin{smallpmatrix} s\circ p&\alpha\circ r \end{smallpmatrix}} \ar[d]^{\begin{smallpmatrix} \id& 0\\0&\beta
\end{smallpmatrix}} & E_0 \ar[d]^{\beta} \ar[r] & 0 \\
0\ar[r] & E_0 \ar[r]_-{\begin{smallpmatrix} j\circ r\\ \beta
\end{smallpmatrix}} & E\oplus E_2 \ar[r]_-{\begin{smallpmatrix} p&0
\end{smallpmatrix}} & E_2 \ar[r] & 0
}
\end{gathered}
\end{equation*}
represents $\beta^*(\alpha_*(\xi))$. We will later use this first line in the form
\begin{equation}
\label{equation:beta(alpha(xi))}
0\to E_0\To{\gamma_1} E_0\oplus E\To{\delta_1} E_0 \to 0,
\end{equation}
with $\gamma_1=\bigl(\begin{smallmatrix} s\circ\beta\\j\circ r\end{smallmatrix}\bigr)$ and $\delta_1=(\begin{matrix} \alpha\circ r&s\circ p \end{matrix})$.

We next look at the restriction of the universal extension over the embedded double point $\nearrow$ at $0$ in $W$, which points in the direction of $\xi$. We will write $2X:=X\times \nearrow\subset X\times W$, $X:=X\times 0\subset X\times \nearrow\subset X\times W$ and denote by $\cO_X[t]:=\cO_X\otimes_\C\C[t]=\cO_{2X}$ the structure ring of $2X$ and by $\pi:2X\to X$ the projection. The class of this extension will be given by $t\pi^*(\xi)\in \Ext^1_{\cO_{2X}}(E_{2,2X},E_{1,2X})$, where $E_{i,2X}:=\pi^*E_i$. The multiplication by $t$ on $\Ext^1_{\cO_{2X}}(F_2,F_1)$ is given by $\mu_*:\Ext^1_{\cO_{2X}}(F_2,F_1)\to\Ext^1_{\cO_{2X}}(F_2,F_1)$, where $\mu=\mu_t:F_1\to F_1$ is the multiplication morphism by $t$ on $F_1$. We apply it to the element $\pi^*(\xi)$ which is represented by the pull-back of the extension \eqref{equation:xi} to $2X$. We~first note that the inverse image $\pi^*F=F\otimes_\C\C[t]$ to $2X$ through $\pi$ of a $\cO_X$\nobreakdash-module~$F$ is isomorphic as a $\cO_X$-module to $F\oplus F$. On $F\oplus F$ multiplication by $t$ is given by the $\cO_X$-linear operator
$\begin{smallpmatrix}0& 0\\ \id&0\end{smallpmatrix}$ which gives $F\oplus F$ its $\cO_X[t]$-module structure back. In~these terms the extension \eqref{equation:xi} pulls back to $2X$ as
\begin{equation}
0\to E_1\oplus E_1 \To{\bigl(\begin{smallmatrix} j&0\\0&j\end{smallmatrix}\bigr)} E\oplus E\To{\bigl(\begin{smallmatrix} p&0\\0&p \end{smallmatrix}\bigr)} E_2\oplus E_2 \to 0
\end{equation}
and the lower line of the following diagram of $\cO_{2X}$-modules represents $t\pi^*(\xi)$:
\begin{equation}
\label{equation: t xi}
\begin{gathered}
\xymatrix@C=1.3cm{
0\ar[r]
& E_1\oplus E_1
\ar[d]_-{\begin{smallpmatrix} 0& 0\\ \id&0 \end{smallpmatrix}} \ar[r]^-{\begin{smallpmatrix} j&0\\0&j\end{smallpmatrix}}
& E\oplus E
\ar[d]^-{\begin{smallpmatrix} 0& s\circ p\\ \id&0
\end{smallpmatrix}} \ar[r]^-{\begin{smallpmatrix} p&0\\0&p \end{smallpmatrix}}
& E_2\oplus E_2 \ar[d]^{\id} \ar[r] & 0
\\
0\ar[r]
& E_1\oplus E_1
\ar[r]_-{\begin{smallpmatrix} \alpha&0\\0&j
\end{smallpmatrix}}
& E_0\oplus E
\ar[r]_-{\begin{smallpmatrix} 0&p\\ \beta&0
\end{smallpmatrix}}
& E_2\oplus E_2 \ar[r] & 0,
}
\end{gathered}
\end{equation}
where the $t$-multiplication on the term $E_0\oplus E $ is given by the operator $\begin{smallpmatrix} 0& s\circ p\\ j\circ r&0
\end{smallpmatrix}$. We~may write this line also in the form
$ 0\to \pi^*E_1\to\cE\to\pi^*E_2\to0$. Tensoring it by $0\to\cO_X\To{t}\cO_{2X}\to\cO_X\to0$ leads to a commutative diagram
\begin{equation}
\label{equation: diagram 3x3}
\begin{gathered}
\xymatrix@R=.5cm{
&0\ar[d]&0\ar[d]&0\ar[d]&
\\
0\ar[r]
& E_1 \ar[d]_{\alpha} \ar[r]
& \pi^*E_1\ar[d] \ar[r]
& E_1 \ar[d]^{\alpha} \ar[r] & 0
\\
0\ar[r]
& E_0\ar[d]_\beta\ar[r]^-{\gamma_2}
&\cE\ar[d]\ar[r]^-{\delta_2}
& E_0\ar[d]^\beta\ar[r] & 0
\\
0\ar[r]
& E_2 \ar[d]_{\alpha} \ar[r]
& \pi^*E_2 \ar[r]\ar[d]
& E_2 \ar[d] \ar[r] & 0
\\
&0&0&0&
}
\end{gathered}
\end{equation}
of $\cO_{2X}$-modules with exact rows and columns and the extension class of its middle row is the first order deformation induced by the family $\cE$, see \cite[Th.\,2.7]{Hartshorne_Deformations}. If we replace now the morphisms in the middle row by those of the sequence \eqref{equation:beta(alpha(xi))} representing $\beta^*(\alpha_*(\xi))$ we get again a diagram of $\cO_{2X}$-modules with exact rows and columns, \cf \eqref{equation: t xi},
\begin{equation}
\label{equation: diagram 3x3 new}
\begin{gathered}
\xymatrix@R=.35cm{
&0\ar[d]&0\ar[d]&0\ar[d]&
\\
0\ar[r]
& E_1 \ar[dd]^{\alpha}
\ar[r]^-{\begin{smallpmatrix} 0\\ \id
\end{smallpmatrix}}
& E_1\oplus E_1
\ar[dd]^{\begin{smallpmatrix} \alpha&0\\0&j
\end{smallpmatrix}}
\ar[r]^-{\begin{smallpmatrix} \id& 0
\end{smallpmatrix}}
& E_1 \ar[dd]^{\alpha} \ar[r] & 0
\\ \\
0\ar[r]
& E_0\ar[dd]^\beta\ar[r]^-{\gamma_1}
&E_0\oplus E
\ar[dd]^{\begin{smallpmatrix} 0&p\\ \beta&0
\end{smallpmatrix}}
\ar[r]^-{\delta_1}
& E_0\ar[dd]^\beta\ar[r] & 0
\\ \\
0\ar[r]
& E_2
\ar[d]
\ar[r]_-{\begin{smallpmatrix} 0\\ \id
\end{smallpmatrix}}
& E_2\oplus E_2
\ar[r]_-{\begin{smallpmatrix} \id& 0
\end{smallpmatrix}}
\ar[d]
& E_2 \ar[d] \ar[r] & 0.
\\
&0&0&0&
}
\end{gathered}
\end{equation}
We will show that the two middle rows of \eqref{equation: diagram 3x3} and \eqref{equation: diagram 3x3 new} are equivalent as extensions. In this direction, first notice that the differences $\delta_2-\delta_1$ and $\gamma_2-\gamma_1$ factorise through morphisms of $\cO_{2X}$-modules $E_2\oplus E_2\to E_1$ and $E_2\to E_1\oplus E_1$, respectively. Using the $\mathcal{O}_{2X}$-module structure given by multiplication with $t$, we see that the induced morphisms have the form $\begin{smallpmatrix} u&0
\end{smallpmatrix}$ and $\begin{smallpmatrix} 0\\ v
\end{smallpmatrix}$, respectively, with $u,v\in\Hom_{\cO_X}(E_2,E_1)$. Putting $\epsilon=\begin{smallpmatrix} 0&\alpha\circ u\circ p\\ 0&0
\end{smallpmatrix}$
and
$\epsilon'=\begin{smallpmatrix} 0&0\\ j\circ v\circ\beta&0
\end{smallpmatrix}$ in $\End_{\cO_X}(E_0\oplus E)$
we get $\epsilon^2=(\epsilon')^2=\epsilon\circ\epsilon'=
\epsilon'\circ\epsilon=0$, $\delta_2-\delta_1=\delta_1\circ\epsilon$, $\gamma_2-\gamma_1=\epsilon'\circ\gamma_1$, $\epsilon\circ\gamma_1=0$ and $\delta_1\circ\epsilon'=0$. Hence, the
diagram
\begin{equation*}
\begin{gathered}
\xymatrix@R=1cm{
0\ar[r] & E_0 \ar[d]_{\id}\ar[r]^-{\gamma_2}
& E_0\oplus E\ar[r]^-{\delta_2} \ar[d]^{\id+\epsilon-\epsilon'} & E_0 \ar[d]^{\id} \ar[r] & 0 \\
0\ar[r] & E_0 \ar[r]^-{\gamma_1}
& E_0\oplus E\ar[r]^-{\delta_1} & E_0 \ar[r] & 0
}
\end{gathered}
\end{equation*}
is commutative and shows that the two extensions under consideration lie in the same class in $\Ext^1_{\cO_X}(E_0,E_0)$, which was to be shown.
\end{proof}

\subsection{Semistable coherent sheaves}\label{subsect:coherentsheaves}

We will work on a fixed complex projective manifold $X$ endowed with a cohomology class $\omega \in H^{1,1}(X, \mathbb{R})$ that can be represented by a Kähler form. This class will serve as a polarisation which will help us to define Gieseker-Maruyama-semistability for coherent sheaves on $X$, \cf \cite[Def.\,11.1]{GRT1}. We start by studying basic properties of semistable sheaves. For simplicity, we will only consider the case of torsion-free sheaves, although most properties are valid for pure coherent sheaves. Later on, we will focus on the case of rank two torsion-free sheaves.

\begin{Defi}
Let $E$ be a coherent sheaf on $X$. Its \emph{Hilbert-polynomial (with respect to $\omega$)} is defined as the polynomial function (with coefficients in $\mathbb{C}$) that is given by
\[
P_E(m):=P^{\omega}_E(m) := \int_X \ch(E)\, e^{m\omega}\, \Todd(X),
\]
where $\ch(E)$ and $\Todd(X)$ denote the Chern character of $E$ and the Todd class of $X$, respectively.
If $E$ is torsion-free and non-zero we define its \emph{reduced Hilbert-polynomial} as
\[
p_E:=p^{\omega}_E := \frac{P_E}{\rank E}.
\]
We will say that $E$ is \emph{(Gieseker-Maruyama-)stable (with respect to $\omega$)} and \emph{semistable}, respectively, if $E$ is torsion-free and if for any coherent subsheaf $0\neq F\subsetneq E$ one has $p_F<p_E$ and $p_F\le p_E$, respectively. We will call $E$ \emph{polystable} if it splits as a direct sum of stable subsheaves having the same reduced Hilbert-polynomial.
If $E$ is semistable but not stable we will say that it is \emph{properly semistable}.
\end{Defi}

The usual relations to slope-stability (with respect to $\omega$), which will also be referred to as $\mu$-stability, continue to hold in this context, namely:
\[
\text{$\mu$-stable $\implies$ stable $\implies $ semistable
$\implies $ $\mu$-semistable,}
\]
\cf \cite[Lem.\,1.6.4]{HL}.
In particular, the boundedness result \cite[Prop.\,6.3]{GrebToma} for $\mu$-semistable sheaves implies:

\begin{Prop}[Boundedness]\label{boundedness}
Let $X$ be a $d$-dimensional projective manifold and let $K$ a compact subset of the Kähler cone $\mathcal{K}(X) \subset H^{1,1}(X, \mathbb{R})$ of $X$. Fix a natural number $r>0$ and classes $c_i\in H^{2i}(X,\R)$, $i = 1, \dots, d$. Then, the family
of rank~$r$ torsion-free sheaves $E$ with $c_i(E)=c_i$ that are semistable with respect to some polarisation contained in $K$ is bounded.
\end{Prop}

The proofs of the following three basic results are standard and therefore left to the reader, \cf \cite[Prop.\,1.2.7]{HL}, \cite[Prop.\,3.1]{Seshadri}\footnote{Seshadri formulates and proves the corresponding result for slope-semistable vector bundles of degree zero over a Riemann surface; Gieseker-Maruyama-semistability is the correct higher-dimensional semistability condition to make this work in general.} and \cite[\S 9.3]{LePotier_Lectures}, and finally \cite[Prop.\,1.5.2]{HL}, respectively.

\begin{lemma}\label{lem:morphisms}
Let $E$ and $E'$ be semistable sheaves on $X$ and let $\phi :E\to E' $ be a non-zero morphism of $\cO_X$-modules. Then $p_E\le p_{E'}$. If equality holds, then $\im(\phi)$ is semistable and $p_{\im(\phi)}=p_E= p_{E'}$. If moreover the rank of $\im(\phi)$ coincides with the rank of $E$ or with the rank of $E'$ then $\im(\phi)$ is isomorphic to $E$ or to $E'$ respectively.
\end{lemma}

\begin{Prop}
\label{prop:abelian_category}
The full subcategory $\Coh^{{\rss}}{(X,\omega,p)}$ of the category of coherent sheaves on $X$, whose objects are the semistable sheaves with fixed reduced Hilbert polynomial $p$ and the zero-sheaf, is abelian, noetherian and artinian.
\end{Prop}

\begin{Prop}[Jordan-Hölder filtrations]\label{prop:JH}
Any semistable sheaf on $X$ admits a Jordan-Hölder filtration in the sense of \cite[Def.\,1.5.1]{HL} (with respect to $\omega$-stability). The associated graded sheaf is unique up to isomorphism.
\end{Prop}

The derivation of the following result is less formal and requires deeper insight into the geometry of Douady spaces.

\begin{thm}[Openness of (semi)stability]\label{thm:openness-semistability}
Let $(S,0)$ be a complex space germ and~$\cE$ be a coherent sheaf on $X\times S$ that is flat over $S$. If the fibre of $\cE$ over $0\in S$ is (semi)stable, then the fibres of $\cE$ over any point in a neighbourhood of $0$ in $S$ are likewise (semi)stable.
\end{thm}

\begin{proof}
The proof of \cite[Cor.\,5.3]{TomaLimitareaI} immediately adapts to our situation to show that the relative Douady space $D_S(\cE)_{\le b}$ of quotients of $\cE$ with
degrees bounded from above by $b$ is proper over $S$; see \cite{TomaLimitareaII} for details. Using this as well as \cite[Lem.\,4.3]{TomaLimitareaI} to replace Grothendieck's Lemma, we may then prove openness of (semi)stability as in the classical case of ample polarisations, as presented for example in~\cite[Prop.\,2.3.1]{HL}.
\end{proof}

\subsection{Semistable sheaves of rank two}\label{subsect:rank2}

The next result gives a classification of semistable sheaves of rank two on a fixed projective manifold $X$ that is endowed with a given Kähler form $\omega$ and computes the automorphism group for all the resulting classes.

\begin{Prop}[Classification of semistable sheaves]\label{prop:ss sheaves}
Any semistable sheaf $E$ of rank $2$ on $X$ falls into exactly one of the following classes:
\begin{enumerate}
\item\label{prop:ss sheaves1} \emph{Polystable sheaves}
\begin{enumerate}
\item\label{prop:ss sheaves1a}
{\em Stable sheaves.} In this case $\Aut(E)\cong\C^*$.
\item\label{prop:ss sheaves1b}
{\em Decomposable sheaves of the form $L_1\oplus L_2$ with $L_1\ncong L_2$ and $P_{L_1}=P_{L_2}$}. In this case $\Aut(E)\cong\C^*\times \C^*$, and $\Hom(L_1,E)$, $\Hom(E, L_1)$, $\Hom(E,L_2)$, as well as $\Hom(L_2,E)$ are one-dimensional.
\item\label{prop:ss sheaves1c}
{\em Decomposable sheaves of the form $L\oplus L$}. In this case $\Aut(E)\cong\GL(2,\C)$ and $\Hom(L,E)$, $\Hom(E, L)$ are two-dimensional.
\end{enumerate}
\item\label{prop:ss sheaves2}
\emph{Non-polystable sheaves}
\begin{enumerate}
\item\label{prop:ss sheaves2a}
{\em Centres of non-trivial extensions of the form $0\to L_1\to E\to L_2\to 0$ with $L_1\ncong L_2$ and $P_{L_1}=P_{L_2}$}. In this case, we have $\Aut(E)\cong\C^*$, $\Hom(L_1,E)\cong\C$, $\Hom(E, L_1)=0$, $\Hom(E,L_2)\cong\C$, and $\Hom(L_2,E)=0$.
\item\label{prop:ss sheaves2b}
{\em Centres of non-trivial extensions of the form $0\to L\To{\alpha} E\To{\beta} L\to0$.} In this case $\Aut(E)=\{ a\cdot \Id_E+c\cdot\alpha\circ\beta \ | \ a\in\C^*, \ c\in\C\}\cong\C^*\times\C$, $\Hom(L,E)\cong\C$ and $\Hom(E,L)\cong\C$.
\end{enumerate}
\end{enumerate}
In all cases listed above, $L_1$, $L_2$, $L$ are torsion-free sheaves of rank one on $X$.
\end{Prop}
\begin{proof}
The classification follows easily from the existence and uniqueness of Jordan-Hölder filtrations, see Proposition~\ref{prop:JH}. We will hence only compute the automorphism groups and the homomorphism groups here, relying mostly on Lemma~\ref{lem:morphisms}. The three cases listed under \eqref{prop:ss sheaves1} are clear. To deal with the cases listed under \eqref{prop:ss sheaves2}, let~$E$ be the centre of a non-trivial extension of the form
\[
0\to L_1\To{\alpha} E\To{\beta} L_2\to0
\]
with $P_{L_1}=P_{L_2}$.

In case $L_1\ncong L_2$, using the fact that the extension is assumed to be non-split we immediately get $\Hom(L_1,E)\cong\C$, $\Hom(E, L_1)=0$, $\Hom(E,L_2)\cong\C$, and $\Hom(L_2,E)=0$. Applying now $\Hom(E,\cdot)$ to the defining exact sequence of $E$ we obtain $\Hom(E,E)\cong\C$, hence $\Aut(E)\cong\C^*$.

Suppose now that $L_1\cong L_2=:L$. Let $\sigma\in\Hom(E,L)$. Then, $\sigma\circ\alpha=0$, otherwise~$\sigma$ would be a retraction of $\alpha$, contradicting the assumption that the extension is non-split. Consequently, $\sigma$ factors through $\beta$, \ie $\sigma=c\beta$ for some $c\in\C$. In particular, $\Hom(E, L)\cong\C$. Similarly we get $\Hom(L,E)\cong\C$. Applying as before $\Hom(E,\cdot)$ to the defining exact sequence of $E$, we get
\[
0\to\Hom(E,L)\To{\alpha\circ \cdot}\Hom(E, E)\To{\beta\circ\cdot}\Hom(E, L).
\]
The image of an element $\phi$ in $\Hom(E, E)$ through
the map $\Hom(E, E)\To{\beta\circ\cdot}\Hom(E, L)$ is of the form $a\beta$ for some
$a\in\C$, with $a\neq0$ if $\phi\in\Aut(E)$.
With this notation $\beta\circ(\phi-a\Id_E)=0$, hence $\phi-a\Id_E=\alpha\circ c\beta=c\cdot \alpha\circ\beta$ and the desired description of $\Aut(E)$ follows.
\end{proof}

\begin{cor}\label{cor:extensions}
Up to a multiplicative constant every
semistable non-polystable sheaf $E$ of rank $2$ on $X$ gives rise to a unique extension
\[
0\to L_1\to E\to L_2\to 0,
\]
with rank one torsion free sheaves $L_1$, $L_2$ on $X$ such that $P_{L_1}=P_{L_2}$.
\end{cor}

\subsection{Basic geometric properties of the stack of semistable sheaves} \label{subsect:quotient_presentations}

In our basic conventions regarding \emph{algebraic spaces} we follow \cite[\S2]{Alper13}. In particular, algebraic spaces are sheaves in the étale topology and their diagonal morphisms are assumed to be quasi-compact. This is in line with the conventions adopted in Knutson's book \cite{Knutson}. We point out that the conventions of the Stacks Project are a priori different, see \cite[\href{https://stacks.math.columbia.edu/tag/025X}{Tag\,025X}]{stacks-project}, but eventually equivalent, see \cite[\href{https://stacks.math.columbia.edu/tag/076L}{Tag\,076L}]{stacks-project}. In our basic conventions regarding \emph{algebraic stacks} we follow \cite{stacks-project}.

We consider the stack $\cX:=\cC\!oh_{(X,\omega),\tau}^{{\rss}}$ of semistable sheaves on $(X,\omega)$ with fixed rank and Chern classes; the latter data will be collected in a vector $\tau=(r, c_1, \dots, c_{2\dim X})$, which we call the \emph{type} of the sheaves under consideration. This is an algebraic stack locally of finite type over $\C$ since it satisfies Artin's axioms \cite[Th.\,2.20]{Alper}; see also \cite[Th.\,2.19]{AlperHR}.

\subsubsection{Quotient stack realisation}\label{subsubsect:quotient_realisation}

The stack $\cX$ may be realised as a quotient stack in the sense of \cite[Def.\,3.1]{Alper} in the usual way; we quickly recall the construction, which is explained for example in \cite[\S 4.3]{HL} or \cite[Ex.\,8.7]{Alper13}: choose an ample line bundle $\cO_X(1)$ and an integer $m$ such that all semistable sheaves (with respect to~$\omega$) with fixed rank and Chern classes $\tau$ on $X$ are $m$-regular with respect to $\cO_X(1)$.
This is possible since we have boundedness for such sheaves by Proposition \ref{boundedness}. Since the rank and the Chern classes of the sheaves $F$ under consideration are fixed, by $m$-regularity and Riemann-Roch we obtain that $h^0(F(m))$ is constant, equal to $N \in \mathbb{N}$.
Setting $V:=\C^N$, $\cH:=V\otimes_\C\cO_X(-m)$, we obtain for any $F$ as above an epimorphism of $\cO_X$-modules
$\rho:\cH\to F$
as soon as we have fixed an isomorphism $V\to H^0(F(m))$.
Moreover, the induced map $H^0(\rho(m)):H^0(\cH(m))\to H^0(F(m))$ is bijective. We thus get a point $[\rho:\cH\to F]$ in the open (quasi-projective) subscheme $R$ of $\Quot_\cH$ of semistable quotients $F$ of $\cH$ with type $\tau$ that induce isomorphisms at the level of $H^0(\rho(m)):H^0(\cH(m))\to H^0(F(m))$. The natural action of the linear group $G:=\GL(V)$ on $V$ induces an action on $\Quot_\cH$ leaving the open subset $R$ invariant. Let $\cF$ be the universal quotient sheaf restricted to $X\times R$. It is a $G$-sheaf and it allows to define an isomorphism from the quotient stack $[R/G]$ to $\cX$. Indeed, an object of $[R/G]$
is a triple $(T,\pi:P\to T, f:P\to R)$, where $T$ is a scheme, $\pi$ is a principal $G$-bundle and $f$ is a $G$-equivariant morphism. Then the $G$-sheaf obtained from $\cF$ by pullback to $X\times P$ gives a flat family of semistable sheaves on $X$ parametrised by $T$ and thus an object of $\cX$.
Conversely if $\cE$ is a flat family of semistable sheaves of type $\tau$ on~$X$ parametrised by a scheme $S$, then as in the proof of \cite[Lem.\,4.3.1]{HL} the frame bundle $R(\cE(m))$ associated to it gives an object $(S,R(\cE(m))\to S,R(\cE(m))\to R)$ of $[R/G]$.

As one consequence of the above analysis, we notice that $\cX$ has affine diagonal morphism; in particular, $\cX$ is
an \emph{Artin stack} in the sense of \cite[\S2]{Alper13}. Indeed, using the chart $R\to\cX$, in order to check that the map $\cX\to\cX\times_{\Spec \C}\cX$ is affine comes down to showing that $G\times R\to R\times R$, $(g,q)\mto(q,g\acts q)$ is affine. But this is a map of $R$-schemes, where $G\times R$ is affine over $R$ and $R\times R$ is separated over $R$, hence the conclusion follows by part (2) of \cite[\href{https://stacks.math.columbia.edu/tag/01SG}{Tag\,01SG}]{stacks-project}.

\subsubsection{Closed points and closures of points} We will characterise closed points in terms of polystability and show that polystable degenerations are unique. Grauert semicontinuity, see \cite[Prop.\,III.12.8]{Hartshorne}, is the key principle at work here. With a view towards the discussion carried out in subsequent parts of the paper we will restrict ourselves to the case of coherent sheaves having rank two.

In the subsequent discussion, we will use the following notation: If $G$ is an algebraic group and $X$ is a $G$-scheme, then for $x \in X(\mathbb{C})$ we denote by $[x]_G$ the image of $x$ under the morphism $X \to [X/G]$. We will also use the same notion for the associated points in the corresponding topological spaces $|X|$ and $|[X / G]|$.

\begin{Prop}[Characterising closed points]\label{prop:characterising_closed_points}
Let $z \in R$ be a closed point. Then, the following are equivalent.
\begin{enumerate}
\item\label{prop:characterising_closed_points1} The point $[z]_G \in |[R/G]| \cong |\cX|$ is closed.
\item\label{prop:characterising_closed_points2} The $G$-orbit $G\acts z \subset R$ is closed.
\item\label{prop:characterising_closed_points3} The sheaf $\cF_z$ is polystable.
\end{enumerate}
\end{Prop}
\begin{proof}
The equivalence ``\eqref{prop:characterising_closed_points1} $\Leftrightarrow$ \eqref{prop:characterising_closed_points2}'' follows directly from the definitions. In order to show ``\eqref{prop:characterising_closed_points2}/\eqref{prop:characterising_closed_points1} $\Rightarrow$ \eqref{prop:characterising_closed_points3}'', assume that $\cF_z$ is non-polystable. Then, $\cF_z$ can be realised as a non-trivial extension $0 \to L_1 \to \cF_z \to L_2 \to 0$, see Proposition~\ref{prop:ss sheaves}. Consequently, $\cF_z$ degenerates to $L_1 \oplus L_2$ over the affine line, and therefore does not give a closed point of $\cX$. It remains to show that orbits of polystable sheaves are closed. This however follows as in \cite[Lem.\,4.7]{Gieseker}, whose proof we reproduce here: Suppose that~$\cF_z$ is polystable and and that $F$ is a flat family of semistable sheaves over a non-singular curve $C$ such that for some point $P\in C$ one has $F_Q\cong\cF_z$ for all $Q\in C\setminus\{ P\}$. Let $E$ be any stable sheaf on $X$ with $p^\omega_E=p^\omega_{\cF_z}$. By upper semicontinuity of the function $Q\mto \Hom(E, F_Q)$ we see that $F_P$ contains at least as many copies of $E$ as are contained in $\cF_z$. Thus $F_P\cong \cF_z$.
\end{proof}
The following now is a consequence of Proposition~\ref{prop:ss sheaves}.
\begin{cor}\label{cor:reductive_stabiliser}
Every closed point of $\cX$ has linearly reductive stabiliser.
\end{cor}
Next, we look at closures of non-closed points.

\begin{Prop}[Uniqueness of polystable degenerations]
For any $\C$-point $y$ of~$\cX$, there exists a unique closed point in $\overline{\{y\}} \subset |\cX|$.
\end{Prop}
\begin{proof}
If $\cF_z$ is polystable, by Proposition~\ref{prop:characterising_closed_points} the corresponding point $y = [z]_G$ is closed, so there is nothing to show.

If $\cF_z$ is semistable but not polystable, then clearly the closed point corresponding to the polystable sheaf $\mathrm{gr}^{JH}(\cF_z)$ lies in $\overline{\{[z]_G\}}$.
Suppose that there is another closed point $x$ of $|\cX|$ lying in $\overline{\{[z]_G\}}$, and let $\cE$ be a polystable sheaf representing $x$. Let
\[0\to L_1\to \cF_z \to L_2\to 0\]
be the unique extension
with rank one torsion free sheaves $L_1$, $L_2$ on $X$ such that $P_{L_1}=P_{L_2}$ as ascertained by Corollary
\ref{cor:extensions}.
Grauert semicontinuity then implies
\begin{align*}
\dim_\C \Hom(\cE,\cF_z) &\geq \dim_\C\Hom(\cF_z, \cF_z),\\
\dim_\C \Hom(L_1, \cE) &\geq \dim_\C\Hom(L_1, \cF_z), \, \text{ and}\\
\dim_\C \Hom(\cE,L_2) &\geq \dim_\C\Hom(\cF_z, L_2),
\end{align*}
from which together with Proposition \ref{prop:ss sheaves} we deduce that the polystable sheaf $\cE$ has to be isomorphic to $ L_1 \oplus L_2 \cong \mathrm{gr}^{JH}(\cF_z)$.
\end{proof}

\subsubsection{Slices and local quotient presentations}
We note that by construction $R$ admits $G$-equivariant locally closed embeddings into the projective spaces associated with finite-dimensional complex $G$-representations, arising from natural $G$-linearised ample line bundles on the Quot-scheme induced by $\cO_X(1)$, see \cite[p.~101]{HL}. This fact will be used in the proof of the subsequent result, which provides rather explicit local quotient presentations for the stack $\cX$. We continue to use the notation established in Section~\ref{subsubsect:quotient_realisation}.

\begin{Prop}[Local quotient presentation induced by slice]\label{prop:slice}
Let $E$ be a semi\-stable sheaf on $X$ corresponding to a closed point $x \in \cX (\mathbb{C})$. Let $s \in R$ project to the closed point $[s]_{G} \in [R/G]$ that is mapped to $x$ by the isomorphism $[R/G] \cong \cX$ established above. Then, there exists a $G_s$-invariant, locally closed, affine subscheme $S$ in $R$ with $s \in S$ such that $T_sR = T_sS \oplus T_s(G\cdot s)$, such that the morphism $G \times S \to R$ is smooth, and such that the induced morphism $f: [S/G_s] \to \cX$ is étale and affine, maps the point $0 :=[s]_{G_s} \in [S/G_s] (\mathbb{C})$ to $x$, and induces an isomorphism of stabiliser groups $G_s = \mathrm{Aut}_{[S/G_s]}(0) \overset{\cong}{\longrightarrow}\mathrm{Aut}_{\cX}(x) \cong \mathrm{Aut}(E)$; \ie $f$ is a local quotient presentation of $\cX$ at $x$ in the sense of \cite[Def.\,2.1]{AlperFS}.
\end{Prop}
\begin{proof}
By Corollary~\ref{cor:reductive_stabiliser}, the stabiliser subgroup $\Aut_{\cX}(x) \cong \Aut(E)$ is linearly reductive. Consequently, the proof of the claim presented in Remark~3.7 and Lemma~3.6 of \cite{AlperKresch} continues to work even without the normality assumption made there, if we replace the application of Sumihiro's Theorem (which uses the normality assumption) by the observation made in the paragraph preceding the proposition that in our setup right from the start $R$ comes equipped with a $G$-equivariant locally closed embedding into the projective space associated with a finite-dimensional complex $G$-representation. Alternatively, see \cite[Props.~9.6 and 9.7]{JoyceSong}.
\end{proof}

\begin{cor}[Slice is stabiliser-preserving]\label{cor:stab_preserving}
In the setup of Proposition~\ref{prop:slice}, let $t \in S \subset R$ and let $H_t$ be the stabiliser group of the action of $H := G_s$ on $S$ at the point $t\in S$. Then, we have $H_t = G_t$.
As a consequence, we obtain
\[H_t \cong \mathrm{Aut}(\mathcal{F}_t)\]
under the morphism of stabiliser groups induced by $f: [S/G_s] \to \cX$.
\end{cor}
\begin{proof}
As the $H$-action on $S \subset R$ is obtained by restricting the $G$-action to the subgroup $H$, we clearly have the inclusion
\begin{equation}\label{eq:one_inclusion}
H_t \subset G_t \quad \quad \forall t \in S.
\end{equation}
Moreover, as for all $t \in S$ the stabiliser subgroup $G_t$ is isomorphic to the automorphism group $\mathrm{Aut}(\mathcal{F}_t)$ of the corresponding member of the family $\mathcal{F}$, and is therefore connected by Proposition~\ref{prop:ss sheaves}, it suffices to show that the two groups appearing in \eqref{eq:one_inclusion} have the same dimension. Consider the twisted product $G\times_H S$, which is the quotient of $G \times S$ with respect to the proper $H$-action given by $h\acts (g, t) := (gh^{-1}, h\acts t)$. The fact that $f: [S/G_s] \to \cX$ is étale implies that the natural $G$-equivariant morphism
\[\varphi: G \times_H S \to R, [g, t]\; \mto g \acts t\]
is étale. In particular, the restriction of $\varphi$ to any $G$-orbit is étale. We conclude that
\[\dim H_t = \dim G_{[e,t]} = \dim G_{\varphi([e,t])} = \dim G_t,\]
as desired.
\end{proof}

The subsequent results will be crucial for the proof of our main result, as it relates abstract deformation theory to the concrete group actions appearing in our setup, see Section~\ref{subsubsect:main_result_main_case}.

\begin{Prop}[Slice provides semi-universal deformation]\label{prop:slice-deformation}
In the situation of Propositions~\ref{prop:slice}, the analytic germ $(S^{\an}, s)$ of $S^{\an}$ at $s$ together with the restriction $(\cU^{\an}, s) := (\cF|_{(S,s) \times X})^{\an}$ of the universal family $\cF$ of $R$ to $(S^{\an}, s)$ is a semi-universal deformation of $E$.
\end{Prop}
\begin{proof}
As both $[S/G_s]$ and $\cX$ are algebraic stacks, there exist formal miniversal deformations $\widehat{\Def}(x)$ and $\widehat{\Def}([s])$ of $x\in \cX(\C)$ and $[s] \in [\Spec A/G_x](\C)$. Moreover, the local quotient presentation establishes an isomorphism of formal schemes $\hat{f}:\widehat{\Def}([s]) \to \widehat{\Def}(x)$. We will check that the first space is isomorphic to the formal completion $\widehat S$ of $S$ at $s$.

We claim that the natural morphism $(S, s)\to ([S/G_s], 0)$ is formally versal at $s$ in the sense of \cite[Def.\,A.8]{AlperHR}. For this, we check the assumptions of \cite[Prop.\,A.9]{AlperHR}: Both $s$ and $0$ are closed points. Moreover, the morphism $S \to [S/G_s]$ is representable and smooth. Hence, the induced map of the $0$-th infinitesimal neighbourhoods $S^{[0]} \to [S/G_s]^{[0]}$ is likewise representable, and for every $n \in \mathbb{N}$ the induced map of $n$-th infinitesimal neighbourhoods $S^{[n]} \to [S/G_s]^{[n]}$ is smooth. Finally, the stabiliser of $[S/G_s]$ at $s$, which is equal to $G_s \cong \mathrm{Aut}(E)$, is reductive. Consequently, part (2) of \cite[Prop.\,A.9]{AlperHR} implies that $(S, s)\to ([S/G_s], 0)$ is formally versal at~$s$, as claimed. Moreover, as $s$ is a $G_s$-fixed point, the induced map on tangent spaces $T_s S \to T_0 [S/G_s]$ is an isomorphism.

As a consequence, we see that the restriction $\widehat{\cU}$ of the universal family $\cF$ to the formal completion $\widehat S$ of $S$ at $s$ is an object of $\cX$ over $\widehat S$ that is formally miniversal at~$s$ in the sense of \cite[Def.\,2.8]{Alper}. Moreover, $\cU^{\an} = (\cF|_{S \times X})^{\an}$ obviously provides an analytification of $\widehat{\cU}$. It follows from the fact that a versal deformation of $E$ exists and from \cite[Satz 8.2]{FlennerHabil} that the germ $(S,s)$ of $S$ at $s$ together with the restriction of $\cU^{\an}$ to this germ is a semi-universal deformation of $E$.
\end{proof}

\begin{rem}
Using analytic stacks, an alternative proof can be given as follows: As in the above proof, one easily checks that the map $(S,s)^{\an} \to [S/G_s]^{\an} = [S^{\an}/G_s]$ is smooth and the induced map on tangent spaces is an isomorphism. These two conditions are equivalent to the conditions in the definition of a semi-universal family, \cf \cite[p.~19]{KosarewStieber90}.
\end{rem}

We also note two properties of the $\mathrm{Aut}(E)$-action on its semi-universal space $(S, s)^{\an}$:

\begin{lemma}[Action of the homothety subgroup]\label{lemma:action_of_homotheties}
In the situation of Proposition~\ref{prop:slice}, the subgroup of homotheties $\mathbb{C}^*\cdot \Id_E$ of $E$ acts trivially on the semi-universal deformation space $S$.
\end{lemma}

\begin{proof}
Under the identification of $\mathrm{Aut}(E)$ with $G_s \subset \GL(V)$, the subgroup $\mathbb{C}^*\cdot \Id_E$ is mapped to $\mathbb{C}^*\cdot \Id_V$, which acts trivially on $\mathrm{Quot}_\cH$, see \cite[proof of Lem.~4.3.2]{HL}.
\end{proof}

\begin{lemma}[Action on tangent space]\label{lem:action_on_tangent}
Using the identification of $\mathrm{Aut}(E)$ with $G_s$, the tangent space of $(S, s)$ is $\Aut(E)$-equivariantly isomorphic to $\E (E,E)$, where the action on the latter space is as described in Section~\ref{subsect:sheafextensions}.
\end{lemma}

\section{Construction of the moduli space}\label{sect:construction}

The aim of this section is to construct a good moduli space for the stack $\cX:=\cC\!oh_{(X,\omega),\tau}^{{\rss}}$ of semistable sheaves with fixed type $\tau = (2, c_1, \dots, c_{2 \dim X})$ on $(X,\omega)$. As we have seen in Section~\ref{subsect:quotient_presentations}, whose notation we will use in our subsequent arguments, this is an algebraic stack with affine diagonal and locally of finite type over $\C$, which can be realised as a quotient stack $\cX \cong [R/G]$. Using this global quotient presentation as well as the local slice models also constructed in Section~\ref{subsect:quotient_presentations}, we will prove that the algebraic stack $\cX$ admits a good moduli space by checking the conditions of \cite[Th.\,1.2]{AlperFS}.

To set the stage, recall the following fundamental definition from \cite{Alper13}, to which the reader is referred for motivating examples (\eg from the theory of Deligne-Mumford stacks and from Geometric Invariant Theory) and for basic properties.
\begin{Defi}
Let $\cX$ be an algebraic stack with affine diagonal. A morphism $\phi: \cX \to X$ to an algebraic space $X$ is a \emph{good moduli space} if the push-forward functor on quasi-coherent sheaves is exact and if $\phi$ induces an isomorphism $\mathscr{O}_X \cong\phi_*\mathscr{O}_{\cX} $.
\end{Defi}

The following is the main result of this section.
\begin{thm}[Existence]\label{thm:existence} The stack $\cC\!oh_{(X,\omega),\tau}^{{\rss}}$ of $\omega$-semistable sheaves with rank two and fixed Chern classes admits a good moduli space $M^{{\rss}}:=M_{(X,\omega),\tau}^{{\rss}}$ with affine diagonal.
\end{thm}

\begin{proof} Recall the criteria for the existence of a good moduli space given in \cite[Th.\,1.2]{AlperFS}:
\begin{enumerate}
\item\label{crit:existence1} For any $\C$-point $y\in\cX(\C)$, the closed substack $\overline{\{y\}} \subset \cX$ admits a good moduli space.
\item\label{crit:existence2} For any closed point $x\in\cX(\C)$ there exists a local quotient presentation $f:\cW\to\cX$ around $x$ in the sense of \cite[Def.\,2.1]{AlperFS} such that
\begin{enumerate}
\item\label{crit:existence2a} $f$ sends closed points to closed points,
\item\label{crit:existence2b} $f$ is stabiliser preserving at closed points of $\cW$.
\end{enumerate}
\end{enumerate}
Following the structure of these conditions, our proof is subdivided into two big steps, establishing Condition \eqref{crit:existence1} and \eqref{crit:existence2}, respectively. The claim regarding affineness of the diagonal of $M^{{\rss}}$ follows from the construction carried out in the proof of \cite[Th.\,1.2]{AlperFS}, which is summarised in \cite[\S 1.1]{AlperFS}.
\subsection{Condition \eqref{crit:existence1}} If $y$ is closed, the condition is easily verified, as the stabiliser group of $y$ is
linearly reductive.
So, suppose that $y$ corresponds to a semistable sheaf~$F$ appearing as an extension
\[
0\to L_1\to F\to L_2\to 0,
\]
where $L_1$, $L_2$ are rank one sheaves with the same Hilbert-polynomial $P$.

To deal with such extensions we consider the stack of flags $\cD rap_X(P,P)$ whose objects over $S$ are sheaves $\cF_1\subset \cF_2$ on $X\times S$ such that $\cF_1$ and $ \cF_2/\cF_1$ are flat over $S$ with fixed Hilbert-polynomials $P$ and $P$ respectively, \cf \cite[\S 2.A.1]{HL}. In fact since the Hilbert polynomial $P^\omega$ is constant in flat families, $\cD rap_X(P,P)$ is a closed and open substack of the stack $\Quot((X\times\cC\!oh(X))/\cC\!oh(X),\cF)$, where $\cF$ is the universal sheaf on $X\times\cC\!oh(X)$.

We note that the forgetful morphism $\phi:\cD rap_X(P,P)\to\cC\!oh(X)$, $(\cF_1,\cF_2)\mto\cF_2$ is representable, \cf \cite[\href{https://stacks.math.columbia.edu/tag/04YY}{Tag\,04YY}]{stacks-project}. We next claim that $\phi$
is proper in the sense of \cite[Définition 3.10.1]{LaumonMB}. Indeed, for any object of $\cC\!oh(X)$, given by a flat family~$\cF$ of coherent sheaves on $X$ parametrised by a scheme $S$ we get a Cartesian diagram
\begin{equation*}
\label{diagr: }
\begin{gathered}
\xymatrix{
S\times_{\cC\!oh(X)}{\cD rap_X(P,P)}\ar[r] \ar[d]_{\phi_S} & \cD rap_X(P, P) \ar[d]^\phi \\
S \ar[r] & \cC\!oh(X)
}
\end{gathered}
\end{equation*}
in which the first vertical map comes from the universal family of quotients of $\cF$ relative to $S$ of Hilbert polynomial $P$. Thus, the morphism $\phi_S$ is the natural map $\Quot_{\cF/S}(P)\to S$, which is proper by \cite[Cor.\,6.1]{TomaLimitareaII}.\footnote{The cited result covers our setup, in which the Hilbert polynomial is considered with respect to a Kähler class $\omega$. If one restricts to $\cC\!oh_{(X,\omega),\tau}^{{\rss}}$ an alternative argument using the classical version for the properness of the Quot scheme combined with Proposition \ref{boundedness} on the boundedness of semistable sheaves should be possible.}
Hence, also $\phi$ is proper. In our case since $X$ is projective it follows in fact that $\phi$ is even projective, but we will not need this fact.

Let now $P=\frac{1}{2}P_E^\omega$, where $E$ is a coherent sheaf on $X$ of type $\tau$.
We consider the substack $\cD rap^{{\rss}}_X(P,P)\subset\cD rap_X(P,P)$ of sheaves $\cF_1\subset \cF_2$ as before such that moreover $\cF_2$ is semistable of Hilbert polynomial $2P$.
Note that in this situation the quotient $\cF_2/\cF_1$ has rank one and no torsion, since otherwise the saturation of $\cF_1$ in $\cF_2$ would contradict the semistability of $\cF_2$.
We thus get a morphism $\cD rap^{{\rss}}_X(P,P)\to\cM{\rs}(X,P)\times\cM{\rs}(X,P)$, $(\cF_1,\cF_2)\mto(\cF_1,\cF_2/\cF_1)$ whose fibre over a closed point $(L_1,L_2)\in\cM{\rs}(X,P)\times\cM{\rs}(X,P)$ is the closed substack
\[\cD rap_X(L_1,L_2)\subset\cD rap^{{\rss}}_X(P,P)\]
of flags $\cF_1\subset \cF_2$ such that the fibres of $\cF_1$ are isomorphic to $L_1$ and the fibres of $\cF_2/\cF_1$ are isomorphic to $L_2$. Here $\cM{\rs}(X,P)$ denotes the moduli space of rank one stable sheaves on $X$ with Hilbert polynomial $P$.

\subsubsection{The case when $L_1\cong L_2$.} We look at points $y\in \cC\!oh(X)(\C)$ corresponding to coherent sheaves $F$ on $X$ that sit in a short exact sequence of the type
\[
0\to L\to F\to L\to0.
\]
Such coherent sheaves $F$ are semistable with respect to any polarisation on $X$, so also with respect to the ample line bundle $\cO_X(1)$. Let $h=c_1(\cO_X(1))$ and consider the open substack
$\cV:=\cC\!oh_{(X,\omega),\tau}^{{\rss}}\cap\cC\!oh_{(X,h),\tau}^{{\rss}}$ of $\cC\!oh_{(X,\omega),\tau}^{{\rss}}$. Then,
$y\in\cV(\C)$.
The stack $\cD rap_X(L,L)$ is proper over both $\cC\!oh_{(X,\omega),\tau}^{{\rss}}$ and $\cC\!oh_{(X,h),\tau}^{{\rss}}$, and its image contains the respective closures of ${y}$ in $\cC\!oh_{(X,\omega),\tau}^{{\rss}}$ and $\cC\!oh_{(X,h),\tau}^{{\rss}}$, which therefore coincide. It is thus enough to show that the closure $\overline{\{y\}}$ in $\cC\!oh_{(X,h),\tau}^{{\rss}}$ admits a good moduli space.
But the latter stack admits a good moduli space itself \cite[Ex.\,8.7]{Alper13}, so the same will hold for the closed substack $\overline{\{y\}}$ by \cite[Lem.\,4.14]{Alper13}.

\subsubsection{The case when $L_1\ncong L_2$.} For later usage in the subsequent argument, we recall the following well known result.

\begin{lemma}\label{lemma: extensions}
If $L_1\ncong L_2$, then
$\cD rap_X(L_1,L_2)\cong[\E(L_2, L_1) /\C^*\times\C^*]$.
\end{lemma}

\begin{proof} We only sketch the line of the argument, which is most likely already present somewhere in the literature on the subject.

We choose to identify $\C^*\times\C^*$ with $\Aut(L_1)\times \Aut(L_2)$. Then the induced action on
$W:=\E(L_2, L_1)$ is given by $w(\theta_1,\theta_2)=\theta_1\theta_2^{-1}w$.
With this convention if $w\in W$ is the class of an extension
$0\to L_1\To{\alpha} F\To{\beta} L_2\to 0$ then $\theta_1\theta_2^{-1}w$ is represented by the second line of the diagram
\begin{equation*}
\label{diagr: extensions}
\begin{gathered}
\xymatrix@C=1.3cm{
0\ar[r]&
L_1\ar[r]^-\alpha \ar[d]_-{\theta_1\Id_{L_1}} &
F\ar[r]^-\beta \ar[d]^{\Id_{F}} &
L_2\ar[r] \ar[d]^{\theta_2\Id_{L_2}}&
0 \\
0\ar[r]& L_1\ar[r]^-{\theta_1^{-1}\alpha} & F\ar[r]^-{\theta_2\beta} & L_2\ar[r]& 0.
}
\end{gathered}
\end{equation*}

An object of $[W/\C^*\times\C^*]$ is a triple $(T, P\To{\pi} T, P\To{f} W)$, where T is a scheme, $ P\To{\pi} T$ is a principal $\C^*\times\C^*$-bundle and $P\To{f} W$ is a $\C^*\times\C^*$-equivariant morphism, \cf \cite[Def.\,3.1]{Alper}. To such an object we associate an object of $\cD rap_X(L_1,L_2)$ in the following way. By \cite[Chap.\,7]{LePotier_Lectures},
\cite[Cor.\,3.4]{Lange} there exists a universal extension
\begin{equation}
\label{diagr: universal}
0\to L_{1,W}\to \cF\to L_{2,W}\to 0
\end{equation}
on $W\times X$ which we pull back to $P\times X$. The action of $\C^*\times\C^*$ on $W$ induces a $\C^*\times\C^*$-linearisation on the sheaves $L_{1,P}$ and $\cF_P$, which thus descend to $T\times X$ and give the desired object in $\cD rap_X(L_1,L_2)$ over $T$, \cite[Th.\,4.2.14]{HL}.

For the converse we use the fact that $M{\rs}(X,P)$ admits local universal families (\cf \cite[App.\,4D.\,VI]{HL}), so for any object $(\cF_1,\cF_2)$ of $\cD rap_X(L_1,L_2)$ over $S$, one has $\cF_1\cong L_{1,S}\otimes\cL_{1,X}$ and $\cF_2/\cF_1\cong L_{2,S}\otimes\cL_{2,X}$, for suitable line bundles $\cL_{1}, \cL_2$ on~$S$.
Let $P_1\to S$, $P_2\to S$ be the $\C^*$-principal bundles associated with the line bundles $\cL_{1}, \cL_2$ on~$S$ and let $P:=P_1\times_SP_2$. On $P\times X$ we get a ``tautological'' extension
\[
0\to L_{1,P}\to \cF_2\to L_{2,P}\to 0,
\]
which is the pullback of the universal extension \eqref{diagr: universal} by means of some equivariant morphism $P\To{f} W$ by \cite[Cor.\,3.4]{Lange} again. The triple $(S,P\to S, P\To{f} W)$ is the corresponding object of $[W/\C^*\times\C^*]$ that we were looking for.
\end{proof}

As $L_1\ncong L_2$ the natural morphism $\cD rap_X(L_1,L_2)\to \cC\!oh_{(X,\omega),\tau}^{{\rss}}$ is proper with finite fibres, hence finite, \cf \cite[\href{https://stacks.math.columbia.edu/tag/0F2N}{Tag\,0F2N}]{stacks-project}, and therefore affine. Let $\cY_{L_1,L_2}\subset \cC\!oh_{(X,\omega),\tau}^{{\rss}}$ be its image. We will use \cite[Prop.\,1.4]{AlperFS} to show that $\cY_{L_1,L_2}$ admits a good moduli space. From this and \cite[Lem.\,4.14]{Alper13} it will follow that the closed substack $\overline{\{y\}}\subset \cY_{L_1,L_2}$ likewise admits a good moduli space.

Lemma \ref{lemma: extensions} states that $\cD rap_X(L_1,L_2)\cong[\E(L_2, L_1)/\C^*\times\C^*]$, so that in particular $\cD rap_X(L_1,L_2)$ admits a separated good moduli space obtained using classical invariant theory.
It remains to check that $\cY_{L_1,L_2}$ is a global quotient stack and admits local quotient presentations.
As a closed substack of the global quotient stack $\cX\cong[R/G]$, $\cY_{L_1,L_2}$ is a global quotient stack as well, \cf \cite[Def.\,3.4]{Alper}.
Moreover,the morphism $\cD rap_X(L_1,L_2)\to\cY_{L_1,L_2}$ is itself a local quotient presentation for $\cY_{L_1,L_2}$, \cite[Def.\,2.1]{AlperFS}. Indeed, the only condition of that definition which we haven't yet checked is the fact that the morphism $\cD rap_X(L_1,L_2)\to\cY_{L_1,L_2}$ is étale. By \cite[\href{https://stacks.math.columbia.edu/tag/04HG}{Tag\,04HG}]{stacks-project} it is enough to show that the morphism $\cD rap_X(L_1,L_2)\to \cC\!oh_{(X,\omega),\tau}^{{\rss}}$ is unramified, since finiteness has already been established.
This follows from the differential study of the Quot scheme, \cite[Prop.\,2.2.7]{HL}, applied to diagrams of the form\vspace*{-3pt}
\begin{equation}
\label{diagr: Drap(L_1,L_2)}
\begin{gathered}
\xymatrix@R=.5cm{
S\times_{\cC\!oh_{(X,\omega),\tau}^{{\rss}}}{\cD rap_X(L_1,L_2)}\ar[r] \ar[d] & \cD rap_X(L_1,L_2) \ar[d] \\
S \ar[r] & \cC\!oh_{(X,\omega),\tau}^{{\rss}}
}
\end{gathered}
\end{equation}
as before and from the fact that $\Hom(L_1,L_2)=0$, as $L_1\ncong L_2$.

\subsection{Condition \eqref{crit:existence2}}

We next turn our attention to condition \eqref{crit:existence2}. Let $x\in\cX(\C)$ be a closed point and $G_x$ its stabiliser. We will do a case by case analysis depending on the type of a representative $E$ of $x$.\enlargethispage{.8\baselineskip}%

\vspace*{-3pt}\vskip0pt
\subsubsection{The stable case.}
The case when $E$ is stable is quickly dealt with. By openness of stability, see Theorem~\ref{thm:openness-semistability}, it suffices to construct a local quotient presentation at~$[E]$ with the desired properties in the open substack $\cC\!oh_{(X,\omega),\tau}^{\rs} \subset \cX$ of stable sheaves on $X$. We consider the corresponding open $G$-invariant subspace $R^{\rs} \subset R$ of stable quotients and choose a point $p \in R^{\rs}$ mapping to $x$ under the natural map $[R/G] \to \cX$. We note as a first point that every $G$-orbit in $R^{\rs}$ is closed in $R$ by Proposition~\ref{prop:characterising_closed_points}, as a second point that for every point $p \in R^{\rs}$ the stabiliser group $G_p$ is isomorphic to $\mathbb{C}^*$ as $\cE_p$ is simple, and as a third point that it follows from Lemma~\ref{lemma:action_of_homotheties} that $G_p$ acts trivially on the slice $S \ni p$ whose existence is guaranteed by Proposition~\ref{prop:slice} and which, shrinking $S$ if necessary, we may assume to be contained in $R^{\rs}$. As every $G$-orbit in $R^{\rs}$ is closed in $R$, condition \eqref{crit:existence2a} is fulfilled for the quotient presentation induced by $S$, whereas condition \eqref{crit:existence2b} is guaranteed to hold by Corollary~\ref{cor:stab_preserving}. This concludes the discussion of the stable case.

\vspace*{-3pt}\vskip0pt
\subsubsection{The case of a polystable point with $\Aut(E)\cong\GL(2,\C)$.}
If $E$ is polystable with $\Aut(E)\cong\GL(2,\C)$, $x$ is a point in the open substack $\mathcal{V} = \cC\!oh_{(X,\omega),\tau}^{{\rss}}\cap\cC\!oh_{(X,h),\tau}^{{\rss}}$ of $\cC\!oh_{(X,\omega),\tau}^{{\rss}}$. Let $R^{h\text{-ss}}$ be the corresponding $G$-invariant open subscheme of $R$ consisting of $h$-semistable quotients and let $R^{{\rss}}$ be the $G$-invariant open subscheme of $R$ consisting of $\omega$-semistable quotients. Both subschemes contain the $G$-orbit $G\acts p$ corresponding to $E$. As $\cC\!oh_{(X,h),\tau}^{{\rss}}$ admits a good moduli space $M^{h\text{-ss}}$ and as $x$ is a closed point of $\cC\!oh_{(X,h),\tau}^{{\rss}}$, there exists an open, $G$-invariant subscheme $\cU \subset R^{h\text{-ss}} \cap R^{{\rss}}$ that contains $x$ and is saturated with respect to the moduli map $R^{h\text{-ss}} \to M^{h\text{-ss}}$. It follows that the restriction of the moduli map to $\cU$ yields a good quotient $\cU \to \cU\hq G \hto M^{h\text{-ss}}$ for the $G$-action on $\cU$. The desired quotient presentation is then produced by an application of Luna's slice theorem, see for example \cite{D_Luna},
at the closed orbit $G\acts p \subset \cU$.

\vspace*{-3pt}\vskip0pt
\subsubsection{The case of a polystable point with $\Aut(E)\cong\C^*\times\C^*$.} \label{subsubsect:main_result_main_case}

We now start giving the main argument, which is slightly involved and therefore divided into several steps.

\vspace*{-3pt}\vskip0pt
\subsubsection*{Setup} Under the assumption, the point $x$ is in the image $\cY$ of $\cD rap^{{\rss}}_X(P,P)$ which is proper over $\cX$. Moreover, inside $\cD rap^{{\rss}}_X(P,P)$ we have a closed substack corresponding to the condition $\cF_1\cong\cF_2/\cF_1$. Let $\cZ$ be the image of this closed substack in $\cY$. The point $x$ lies in the complement of $\cZ$, so we may assume that the image of the local quotient presentation $f$ guaranteed by Proposition~\ref{prop:slice} is contained in the complement of $\cZ$ too. We will use the notation of that proposition throughout the rest of the proof.

\subsubsection*{Basic properties of the action of $\Aut(E)$} Recall the following properties of the action of $\Aut(E)$ on $S = \mathrm{Spec}(A)$:
\begin{enumerate}
\item The subgroup of homotheties $\C^*\cdot \Id_E \subset \mathrm{Aut}(E)$ acts trivially on $S$, see Lemma~\ref{lemma:action_of_homotheties}; the action of $\Aut(E)$ hence factors over an action of $\mathbb{C}^* \cong (\mathbb{C}^* \times \mathbb{C}^*)/ \mathbb{C}^*$ on $S$.
\item The fibres $F$ over the fixed points for the action of $\Aut(E)$ have $\Aut(F)\cong\C^*\times\C^*$ and conversely, if $\mathrm{Aut}(\mathcal{F}_t) = \C^* \times \C^*$ for some $t \in S$, then $t$ is contained in the $\mathrm{Aut}(E)$-fixed point set in $S$, see Corollary~\ref{cor:stab_preserving}.
\end{enumerate}

The proof is completed once we establish the following\enlargethispage{\baselineskip}%

\begin{key*}
After $S$ has been possibly shrunk further, the orbits of non-polystable fibres $F$ by the $\Aut(E)$-action are not closed.
\end{key*}

The proof of this claim will be given in several steps.

\subsubsection*{Reducing to a low dimensional setup} The set of points of $S$ parametrising non-polystable fibres is constructible by semicontinuity arguments, and so is the set of points belonging to closed orbits by Luna's slice theorem. Both sets are invariant under the $\Aut(E)$-action. If the closure of their intersection does not contain $s$, we just shrink $S$ so that it no longer intersects this closure. So, aiming for a contradiction, suppose that the closure does contain $s$. As this closure is an invariant closed subset of $S$, its image in $S\hq \Aut(E) := \mathrm{Spec}(A^{\Aut(E)})$ is a closed subvariety. Next, we will in a two-step procedure select a one-dimensional analytic space germ inside this image that captures the essential features of the situation.

By cutting down we get an irreducible curve $C$ through the image $o$ of $s$ in $S \hq \Aut(E)$ whose general points correspond to non-trivial closed $\Aut(E)$-orbits in $S$ parametrising non-polystable sheaves.
Let $Z\subset \Spec A$ be an irreducible component of the inverse image of $C$ in $\Spec A$ containing such a general orbit. Without loss of generality, the only point of $Z$ corresponding to a polystable sheaf is $s$. Then, $Z$ is an affine surface with a regular $\C^*$-action and good quotient $\pi :Z \twoheadrightarrow C \subset S \hq \Aut(E)$. As $\pi$ realises the equivalence relation
\[ z_1 \sim z_2 \iff \overline{\mathbb{C}^*\acts z_1} \,\cap\, \overline{\mathbb{C}^* \acts z_2} \neq \emptyset,\]
the fibres of $\pi$ are connected.

Let $(B,o)$ be an analytic branch of $(C,o)$, \ie an irreducible component of the analytic space germ $(C,o)$, and let $Y$ be the closure of $Z_B\setminus Z_o$ in $Z_B$. This is an irreducible analytic space over the germ $(B,o)$. It comes equipped with the naturally induced holomorphic $\C^*$-action with quotient $\pi :Y \twoheadrightarrow (B,o)$, which continues to have connected fibres.\footnote{There is a notion of ``semistable quotients'', which is the holomorphic analogue of the notion ``good quotient'', \eg see \cite{HMP}. Since in our case the quotient map arises as the restriction of a good quotient over a germ inside the original quotient space, we do not need to delve deeper into this theory, though.} We may suppose that $B$ was chosen in such a way that $s$ belongs to $Y$. The fibre $Y_o$ over $o\in B$ is affine and the only point of $Y$ corresponding to a polystable sheaf is $s$. We set $Y^\circ:=Y\setminus\{ s\}$.

The following diagram sums up the relation between the geometric objects constructed so far:
\begin{equation}\label{eq:first_diagram}
\begin{gathered}
\xymatrix@=.7cm{
(Y,s) \ar[dr]& (Y_o, s) \ar[l]\ar[d] \ar[r] & (S,s)\ar[d]\\
&(B,o)\ar[r]&(S \hq \Aut(E), o).
}
\end{gathered}
\end{equation}

\subsubsection*{Introducing universal extension spaces} For simplicity we denote again by $\cF$ the universal sheaf on $X\times Y$.
Any fibre $F$ of $\cF$ over a point $y$ of $Y$ appears as the middle term of an extension $0\to L_1\to F\to L_2\to 0$,
where $L_1$ and $L_2$ are stable of fixed Hilbert polynomial $P$ and non-isomorphic.
If $s \neq y \in Y$, this extension is non-trivial and unique up to a multiplicative constant in $\C^*$, see Corollary \ref{cor:extensions}.
Moreover, the sheaves parametrised by points lying in a fibre of $\pi$ over some point of $B\setminus \{o\}$ are all isomorphic, and thus they all correspond to the same extension, again up to multiplication by a non-zero constant.
Let $\cF_s=L_{1,s}\oplus L_{2,s}$. The natural morphism $\Quot_{\cF/Y}(P)\to Y$ is hence one-to-one over $Y^\circ$, whereas its fibre over $s$ has two points corresponding to $L_{1,s}$ and $ L_{2,s}$. Together with the analytic irreducibility of $Y$, these properties imply that only one of the two quotients $L_{1,s}$ and $ L_{2,s}$ over $s$, say $ L_{2,s}$, is in the closure $Y'$ of $\Quot_{\cF/Y}(P)_{Y^\circ}$ in $\Quot_{\cF/Y}(P)$. We denote by $s'$ the point of $Y'$ lying over $s$ and by $Y'_o$ the fibre over $o$ of the composition $Y'\to Y\to B$.

On $X\times Y'$ we have two rank one universal sheaves, the universal kernel and the universal quotient, which we denote by $\cL_1$ and $\cL_2$, respectively. As explained in the proof of Lemma \ref{lemma: extensions} above, the corresponding extension $0\to \cL_1\to \cF\to \cL_2\to 0$ on $X\times Y'$ restricted to $X\times Y'_o$ gives rise to a
$\C^*\times\C^*$-principal bundle $P:=P_1\times_{Y_o}P_2$ and an equivariant morphism $P\to W$, where $W:=\E( L_{2,s},L_{1,s})$ is the space of extensions of $L_{2,s}$ by $L_{1,s}$. Let $s''$ be any point of $P$ lying over $s'\in Y'_o$. At the level of germs of complex analytic spaces, we get the following commutative diagram extending \eqref{eq:first_diagram}:
\begin{equation}
\label{diagr:germs}
\begin{gathered}
\xymatrix@=.7cm{
& (P,s'')\ar[d]\ar[r] & (W,0)\ar[dd]^\chi\\
(Y',s')\ar[d] & (Y'_o,s')\ar[d]\ar[l] & \\
(Y,s) \ar[dr]& (Y_o, s) \ar[l]\ar[d] \ar[r] & (S,s)\ar[d]\\
&(B,o)\ar[r]&(S \hq \Aut(E), o).
}
\end{gathered}
\end{equation}

\noindent\emph{Analysing the action of $\Aut(E)$.} The automorphism group $\Aut(E) \cong \mathbb{C}^* \times \mathbb{C}^*$ acts equivariantly on the induced diagram of the respective tangent spaces of the germs above.
Recall that by Proposition \ref{prop:slice-deformation} the germ $(S,s)$ may be viewed as the base of the semi-universal deformation of $E$. From Lemma~\ref{lem:action_on_tangent} we obtain an equivariant isomorphism $T_s S \overset{\cong}{\longrightarrow} \E(L_{1,s}\oplus L_{2,s},L_{1,s}\oplus L_{2,s})$. The horizontal map in the second row of the induced equivariant diagram
\[\xymatrix{
T_0W \ar[r]^-{d_0\chi} & T_sS \ar[d]^\cong \\
\E(L_{2,s}, L_{1,s}) \ar[u]^\cong \ar[r]& \E(L_{1,s}\oplus L_{2,s},L_{1,s}\oplus L_{2,s})
}
\]
is the one described by Lemma \ref{lemma:extensions-deformations}, and the group action is induced by the action of $\Aut(L_{1,s})\times\Aut(L_{2,s})\cong\C^*\times\C^*$; in particular, note that the diagonal $\C^*\subset\C^*\times\C^*$ operates trivially on both sides, \cf the proof of Lemma \ref{lemma: extensions}. Almost by definition, the weight of the induced $\C^*$-action on $W$ is $+1$ or $-1$ depending on the isomorphism $\C^*\cong(\C^*\times\C^*)/\C^*$ we have chosen. Without loss of generality, we will assume it to be equal to
$+1$.

Choose a $\mathbb{C}^*$-equivariant closed embedding $\psi :S \hto V$ into a finite-dimensional~$\mathbb{C}^*$ representation space $V$. By composing with the translation by the $\mathbb{C}^*$-fixed point $\psi(s)$ if necessary, we may assume that $\psi(s) = 0 \in V$. Let $V=V_+\oplus V_0\oplus V_-$ be the decomposition of $V$ into subspaces according to the sign of the weights of the $\C^*$-action on $V$. From Diagram~\eqref{diagr:germs} and the consideration regarding the weight of the action on $W$ we infer that $\psi$ embeds $Y_o$ into $V_+$.

We claim that this implies that $\psi$ embeds the reduced space $Y_{\mathrm{red}}$ into $V_+\oplus V_0$. Indeed, if not, there would exist a sequence of points $z_n=(z_{n,+},z_{n,0},z_{n,-})$ in $\psi(Y_{\mathrm{red}})\setminus (V_+\oplus V_0)$ converging to $0$ in $V$.
Then, we can find a sequence of elements $\lambda_n\in \C^*$ with $\lim_{n \to \infty} \lambda_n = 0$ such that $\Vert\lambda_n \acts z_{n,-} \Vert=1$ for each $n\in \N$. It follows that $(\lambda_n \acts (z_{n,+},z_{n,0}))_{n\in \N}$ converges to $0$ in $V_+\oplus V_0$ and that a subsequence of $(\lambda_n \acts z_{n})_{n \in \N}$ converges to a point of norm $1$ in $\psi(Y_{\mathrm{red}})\cap V_-$. Such a point would lie in $Y_o$ contradicting the fact that $Y_o$ is mapped to $V_+$ by $\psi$. Thus, $Y_{\mathrm{red}}$ is embedded into $V_+\oplus V_0$, as claimed.

This analysis finally leads to the desired contradiction: the only closed orbits of the $\C^*$-action on $V_+\oplus V_0$ are the fixed points, whereas by construction $Y$ has many positive-dimensional closed orbits.

This concludes the proof of the key technical claim and hence the proof of Theorem~\ref{thm:existence}, establishing the existence of a good moduli space $M_{(X,\omega),\tau}^{{\rss}}$.
\end{proof}

\section{Properties of the moduli space}\label{sect:properties}
In this section we will establish the functorial properties of $M_{(X,\omega),\tau}^{{\rss}}$ and describe its closed points. Moreover, we will show that $M_{(X,\omega),\tau}^{{\rss}}$ is separated and in fact proper.

\subsection{Functorial properties}
Analogously to \cite[\S 4.1]{HL} we consider the functors $\underline{M'}:=\underline{M'}_{(X,\omega),\tau}:(\textit{Sch}/\C)\to (\textit{Sets})$, $\underline{M}:=\underline{M'}/\sim$, where, for a scheme $S$ over $\C$, $\underline{M'}(S)$ is the set of isomorphism classes of flat families over $S$ of semistable sheaves of type $\tau$ on $X$ and two such families $F,E\in\underline{M'}(S)$ are equivalent through $\sim$
if there exists a line bundle $\cL\in\Pic(S)$ such that $E$ is isomorphic to $F\otimes \cL_X$. As~explained in \cite[\S 4.1]{HL}, if an algebraic space $M$ corepresents the functor $\underline{M'}$, then it also corepresents $\underline{M}$ and the other way round. Finally, using \cite[Lem.\,4.3.1]{HL} and \cite[Th.\,4.16(vi)]{Alper13} we get:

\begin{Prop}\label{prop:corepresentaion} $M_{(X,\omega),\tau}^{{\rss}}$ corepresents the functors $\underline{M'}_{(X,\omega),\tau}$ and $\underline{M}_{(X,\omega),\tau}$.
\end{Prop}

\subsection{Closed points of the moduli space}
Recall that two semistable sheaves on $X$ are called $S$-equivalent if their Jordan-Hölder graduations are isomorphic. By \cite[Th.\,4.16(iv)]{Alper13}, \cite[Lem.\,4.1.2]{HL} and our previous considerations we immediately obtain:

\begin{Prop}[Closed points of moduli space]\label{prop:points} The closed points of $M_{(X,\omega),\tau}^{{\rss}}$ correspond precisely to $S$-equivalence classes of semistable sheaves of type $\tau$ on $X$.
\end{Prop}

\subsection{Separatedness}

The aim of the current section is to prove that the constructed moduli space is separated. This will essentially follow from a refinement of Langton's valuative criterion for separatedness, which is only formulated in \cite{Langton} for two semistable sheaves, at least one of which is stable. But whereas Langton's theorem is stated for slope-semistable sheaves, we are working within the \emph{abelian category} of Gieseker-Maruyama-semistable sheaves of fixed Hilbert polynomial with respect to~$\omega$ (and $0$), \cf Proposition \ref{prop:abelian_category}. This fact is essential for the following discussion to go through.

\subsubsection*{Setup.} We consider a discrete valuation ring $A$ over $\C$ with maximal ideal $\mg$ generated by a uniformising parameter $\pi$. We set $K$ the field of fractions of $A$. We denote by $X_K:=X\times\Spec(K)$ the generic fibre and by $X_\C:=X\times\Spec(\C)$ the special fibre of $X_A:=X\times\Spec(A)$ and by $i:X_K\to X_A$, $j:X_\C:=X\to X_A$ the inclusion morphisms. We denote furthermore by $\xi$ and by $\Xi$ the generic points of $X_\C$ and of~$X_K$, respectively. Note that $\cO_{X_A,\xi}$ is a discrete valuation ring with maximal ideal generated by $\pi$.

\subsubsection*{Langton edges}
Let $E_K$ be a
torsion-free sheaf of rank~$r$ on $X_K$ and let $E\subset i_*E_K$ be a coherent torsion-free sheaf of rank~$r$ on $X_A$ such that $i^*E=E_K$.
Then~$E_\xi$ is a rank~$r$ free $\cO_{X_A,\xi}$-submodule of $(E_K)_\Xi$. Conversely, by \cite[Prop.\,6]{Langton} any rank~$r$ free $\cO_{X_A,\xi}$ submodule $M$ of $(E_K)_\Xi$ gives rise to a unique torsion-free coherent subsheaf~$E$ of $i_*E_K$ on $X_A$ such that $i^*E=E_K$, $E_\xi=M$ and $j^*E$ is torsion-free on $X_\C$.
Langton introduces an equivalence relation on such submodules by putting \hbox{$M\sim\pi^nM$} and calls two equivalence classes $[M]$, $[M']$ {\it adjacent} if there exists a direct sum decomposition $M=N\oplus P$ such that $M'=N+\pi M$.
Equivalent modules induce isomorphic extensions of $E_K$ to coherent subsheaves on $X_A$ as in \cite[Prop.\,6]{Langton}.
This is no longer true in general for adjacent classes. We describe what happens in this case in the following

\begin{rem}\label{rem:adjacent classes}
In the above setup suppose that $[M]$ and $[M']$ are adjacent classes of free rank~$r$ submodules $[M]$ and $[M']$ of $(E_K)_\Xi$ and let $E$, $E'$ denote the coherent sheaf extensions of $M$ and $M'$ to $X_A$. We may suppose that $M$ has a basis $(e_1,\dots,e_r)$ over $\cO_{X_A,\xi}$ such that for a suitable $s\in\{1,\dots,r\}$ the module $M'$ admits $(e_1,\dots, e_s,\pi e_{s+1},\dots,\pi e_r)$ as a basis.
Then the inclusion of $\cO_{X_A,\xi}$-modules $M'\subset M$ induces an inclusion of coherent sheaves $E'\subset E$ on $X_A$ which restricts to a morphism $\alpha: E'_\C\to E_\C$ on $X_\C$, whose image is the unique saturated coherent subsheaf $F$ of $E_\C$ such that $F_\xi$ is the $\cO_{X_\C,\xi}$-vector space generated by $\bar e_1,\dots,\bar e_s$, where the elements~$\bar e_j$ are the images of $e_j$ under $M\to (M\otimes \cO_{X_A,\xi}/\pi \cO_{X_A,\xi})$.
In \cite[Prop.\,7]{Langton} it is shown that $F$ is saturated in $E_\C$ and that $E'$ appears as what is called an elementary transformation of $E$ and in particular that $\Ker(\alpha)\cong \coker(\alpha)$.
One gets two exact sequences
\begin{align}
0&\to \Ker(\alpha)\to E'_\C\to \im(\alpha)\to0, \label{eq:sequence1}\\
0&\to\im(\alpha)\to E_\C\to\coker(\alpha)\to0\label{eq:sequence2}
\end{align}
of torsion-free sheaves on $X_\C$.
Langton calls the passage from $[M]$ to $[M']$ an {\it edge} and denotes it by $[M] -[M']$.
\end{rem}

The following is the key technical observation of this section, eventually proving separatedness of the moduli space. Its proof is modelled after \cite[p.\,101--102]{Langton}.

\begin{Prop}\label{prop:separation} Let $E_K$ be a
torsion-free sheaf of rank~$r$ on $X_K$ and let $E_1, E_2\subset i_*E_K$ be two torsion-free coherent sheaves on $X_A$ such that $i^*E_1=i^*E_2=E_K$ and such that $E_{1,\C}:=j^*E_1$ and $E_{2,\C}:=j^*E_2$ are torsion-free on $X_\C$. Then, the following statements hold.
\begin{enumerate}
\item The sheaves $E_{1,\C}$ and $E_{2,\C}$ have the same (reduced) Hilbert polynomial.
\item If $E_1$, $ E_2$ correspond to an edge with induced morphism $\alpha$ as in Remark \ref{rem:adjacent classes}, and if $E_{1,\C}$ and $\im(\alpha)$ are semistable with reduced Hilbert polynomial $p$, then both $\Ker(\alpha)$ and $E_{2,\C}$ are semistable with reduced Hilbert polynomial $p$. Moreover, $\mathrm{gr}^{JH}(E_{1,\C})\cong\mathrm{gr}^{JH}(E_{2,\C})$.
\item If $E_{1,\C}$ and $E_{2,\C}$ are semistable on $X_\C$, then they can be connected by a finite chain of edges $
[M] -[M']$, $[M'] -[M'']$,\dots, $[M^{(m_{r}-1)}] -[M^{(m_{r})}]$ such that all the associated sheaf extensions $E=E_1$, $E'$,\dots, $E^{(m_r)}=E_2$ to $X_A$ are families of semistable sheaves. Moreover,
\[
\mathrm{gr}^{JH}(E_{1,\C})\cong\mathrm{gr}^{JH}(E_{2,\C}).
\]
\end{enumerate}
\end{Prop}
\begin{proof}
We start by noting that $E_1$ and $E_2$ are flat over $\Spec(A)$ since they are torsion free and $A$ is a discrete valuation ring. Moreover, since they coincide over $\Spec(K)$, their restrictions to $X_\C$ have the same Hilbert polynomial, which proves the first assertion.

If $E_1$, $ E_2$ correspond to an edge and if $E_{1,\C}$ and $\im(\alpha)$ are semistable with reduced Hilbert polynomial $p$, then using the exact sequences in Remark \ref{rem:adjacent classes} together with Proposition~\ref{prop:abelian_category} we see that also $\Ker(\alpha)$ and $E_{2,\C}$ are semistable with reduced Hilbert polynomial $p$. We moreover get $\mathrm{gr}^{JH}(E_{1,\C})\cong\mathrm{gr}^{JH}(E_{2,\C})$, thus establishing the second assertion.

We suppose now that $E_{1,\C}$ and $E_{2,\C}$ are semistable on $X_\C$ and consider now $\cO_{X_A,\xi}$-modules $E_{1,\xi}$, $E_{2,\xi}$. We can find a basis $(e_1,\dots,e_r)$ of $E_{1,\xi}$ over $\cO_{X_A,\xi}$ such that $(\pi^{m_1}e_1,\dots,\pi^{m_r}e_r)$ is a basis of $E_{2,\xi}$, where $m_1$,\dots, $m_r$ are suitable integers.
Up~to replacing $E_2$ by $\pi^nE_2$ for some $n$ and up to permuting the $e_i$-s we may suppose that $m_1=0$ and that $m_1\le m_2\le\cdots\le m_r$.
We will construct a sequence of $m_r$ edges
\[
[M] -[M'],\ [M'] -[M''],\dots,\ [M^{(m_{r}-1)}] -[M^{(m_{r})}]
\]
starting at $[M]=[E_{1,\xi}]$ and ending at $ [M^{(m_{r})}]=[E_{2,\xi}]$ in the following way.
If $s$ is such that $m_s=m_1=0$ and $m_{s+1}>m_s$, we set $M':=(e_1,\dots,e_s,\pi e_{s+1},\dots,\pi e_r)$ to be the $\cO_{X_A,\xi}$-submodule of $(E_K)_\Xi$ generated by the elements $e_1,\dots,e_s,\pi e_{s+1},\dots,\pi e_r$.
We~next set
\begin{align*}
M''&:=(e_1,\dots,e_s,\pi^2 e_{s+1},\dots,\pi^2 e_r),\\[-7pt]&\hspace*{2.7mm}\vdots\\[-7pt] M^{(m_{s+1})}&:=(e_1,\dots,e_s,\pi^{m_{s+1}} e_{s+1},\dots,\pi^{m_{s+1}} e_r)
\end{align*}
and continue with
\[
M^{(m_{s+1}+1)}:=(e_1,\dots,e_s,\pi^{m_{s+1}} e_{s+1},\dots,\pi^{m_{s+1}} e_t, \pi^{m_{s+1}+1} e_{t+1},\dots,\pi^{m_{s+1}+1} e_r),
\]
where $t$ is such that $m_t=m_{s+1}$, $m_{t+1}>m_t$ and so on until we reach $E_{2,\xi}$.
For this sequence of edges we denote by $E=E_1$, $E'$,\dots, $E^{(m_r)}=E_2$ the associated sheaf extensions to $X_A$ and by $\alpha_1,\dots, \alpha_{m_r}$ the induced morphisms.
By construction we have $\im(\alpha_1)=\im(\alpha_1\circ\alpha_2\circ\cdots\circ\alpha_{m_r})$.
By Proposition \ref{prop:abelian_category} it follows that $\im(\alpha_1)$ is semistable with reduced Hilbert polynomial $p:=p_{E_{1,\C}}=p_{ E_{2,\C}}$ so by the second assertion of the proposition also $E'_\C$ is semistable with reduced Hilbert polynomial $p$ and $\mathrm{gr}^{JH}(E_\C)\cong\mathrm{gr}^{JH}(E'_\C)$.
Iterating this piece of argument we thus obtain
\[\mathrm{gr}^{JH}(E_{1,\C})\cong\mathrm{gr}^{JH}(E'_\C)\cong\cdots\cong\mathrm{gr}^{JH}(E^{(m_r-1)}_\C)\cong\mathrm{gr}^{JH}(E_{2,\C}). \qedhere\]
\end{proof}

\begin{rem}\label{rem:proof-of-separation}
The proof of Proposition \ref{prop:separation} works for slope semistability and for more general semistability notions as in \cite[Def.\,2.2]{TomaCriteria} if one replaces the Jordan-Hölder graduation in the category $\Coh(X)$ by the Jordan-Hölder graduation in an appropriate quotient category $\Coh_{d,d'}(X)$ as defined in \cite{HL}. The problem is that while for such semistability notions Jordan-Hölder graduations exist in $\Coh(X)$, they are no longer unique, \cite[Prop.\,2.1]{BuTeTo1bis}.
\end{rem}

In a next step, the sheaf-theoretic construction explained above is re-interpreted in terms of stacks.

\begin{rem}\label{rem:AHLH}
In \cite[\S 3.5]{AHLH} the authors introduce a general framework that allows the characterisation of separatedness and properness of good moduli spaces. We retain here only the interpretation of a Langton edge that they give in
\cite[Rem.\,3.36]{AHLH}. In the above setup, denote by $\overline{ST}_A$ the quotient stack \hbox{$[\Spec (A[s,t]/(st-\pi))/\C^*]$}, where $s$ and $t$ have weights $+1$ and $-1$ with respect to the $\C^*$-action. We note that $\overline{ST}_A\to \Spec (A)$ is a good moduli space. The open locus in $\overline{ST}_A$ where $s\neq0$ is isomorphic to\enlargethispage{\baselineskip}%
\[
[\Spec(A[s,t]_s/(t-\pi/s))/\C^*]\cong [\Spec(A[s]_s)/\C^*]\cong \Spec(A).
\]
As similar remark holds for the open locus where $t \neq 0$. We denote by $[0/\C^*]$ the closed substack $[\Spec(\C)/\C^*]$ of $[\Spec(A[s,t]/(st-\pi))/\C^*]$.

In this setup, it is shown in \cite[Rem.\,3.36]{AHLH} that a Langton edge $[M] -[M']$ with associated sheaf extensions $E$ and $E'$ to $X_A$ gives rise to a morphism $h:\overline{ST}_A\to\cC\!oh(X)$ such that $h|_{\{s\neq0\}}$ and $h|_{\{t\neq0\}}$ correspond to $E$ and $E'$, respectively. One can moreover check using \eqref{eq:sequence1} and \eqref{eq:sequence2} and the discussion at the end of \emph{loc.~cit.} that the coherent sheaf on $X_\C$ corresponding to the composition $\Spec(\C)\to \overline{ST}_A\to\cC\!oh(X)$ is isomorphic to $\Ker(\alpha)\oplus\im(\alpha)$, where $\alpha: E'_\C\to E_\C$ is the morphism of sheaves on $X_\C$ described in Remark \ref{rem:adjacent classes}. Thus, if $E_\C$ and $E'_\C$ are moreover supposed to be semistable, Proposition~\ref{prop:abelian_category} implies that the constructed morphism $h$ takes values in the open substack $\cX$ of $\cC\!oh(X)$ corresponding to semistable sheaves on $X$.
\end{rem}

The proof of the following central result uses Proposition \ref{prop:separation} as well as properties of the good moduli spaces $\cX\to M$ and $\overline{ST}_A\to \Spec (A)$.\footnote{Note that in \cite{AHLH} separation criteria for good moduli spaces are given that apply in very general situations. For the convenience of the reader, we decided to extract parts of the relevant arguments from \emph{loc.~cit.} and to combine them with our results established above in order to give a relatively elementary proof and to make explicit which results of \cite{AHLH} are actually needed in our situation.}

\begin{cor}[Separatedness]\label{cor:separation}
The moduli space $M_{(X,\omega),\tau}^{{\rss}}$
is separated.
\end{cor}

Before giving the proof we state and prove two technical results. While we could not find a reference containing these lemmata,\footnote{But see \cite[\S 7.2]{Kempf} and \cite{Ant}.} their statements are probably well known to experts in the area. They are also not stated here in the most general version in which they hold, but with assumptions that both make our subsequent proof work and allow for a relatively quick proof.

\begin{lemma}[Checking universal closedness by special DVRs]\label{rem:valuative criterion closedness}
Let $k$ be an algebraically closed field and let $f:X\to Y$ be a morphism of algebraic spaces of finite type over $k$. Suppose that $f$ is affine. If for every DVR $A$ with residue field $k$ and fraction field $K$ and for every commutative diagram of solid arrows
\begin{equation*}
\begin{gathered}
\xymatrix{
\Spec(K)\ar[r] \ar[d] & X \ar[d]^f \\
\Spec(A) \ar[r] \ar@{.>}[ru] &Y
}
\end{gathered}
\end{equation*}
there exists a dotted arrow making it commutative, then $f$ is universally closed.
\end{lemma}

\begin{proof}
We adapt the proof of \cite[\href{https://stacks.math.columbia.edu/tag/03KA}{Lem.\,03KA}]{stacks-project} to our situation.
By \cite[\href{https://stacks.math.columbia.edu/tag/03IT}{Lem.\,03IT}]{stacks-project} and the argument in the proof of \cite[\href{https://stacks.math.columbia.edu/tag/03IY}{Lem.\,03IY}]{stacks-project} we may suppose that $Y$ is an affine scheme of finite type over $k$.
Since $f$ is affine we are led to checking the statement when $X$ and $Y$ are affine schemes of finite type over $k$.

Let $T \subset |X|$ be closed. To finish the proof, we have to show that the image $f(|T|)$ is closed. For this, it suffices to show that every closed point $y$ contained in the closure $\overline{f(|T|)}$ already belongs to $f(|T|)$.
Suppose by contradiction that $y$ is a closed point in $\overline{f(|T|)}$ not in $f(|T|)$.
As the image $f(|T|)$ is constructible by Chevalley's Theorem, there exists a (reduced, irreducible) subvariety $Z$ of $Y$ such that $U:= Z \cap f(|T|)$ is open in $Z$ and additionally $y \in Z \setminus U$.
By \cite[Lem.\,7.2.1]{Kempf},
there also
exists an affine curve $C \hto Y$ such that $y \in C \setminus U$ and $C\cap U\neq\emptyset$. There exists then a DVR~$A$ dominating a finite extension of the normalisation of $\mathscr{O}_{C, y}$ whose generic point lifts to~$X$. Applying the assumptions of the lemma to $A$, we conclude that $y \in f(|T|)$.
\end{proof}

\begin{lemma}[Checking separatedness by special DVRs]\label{rem:valuative criteria}
Let $k$ be an algebraically closed field and let $X$ be an algebraic space {\it with affine diagonal and} of finite type over $k$. If for every DVR $A$ with residue field $k$ and fraction field $K$ and for every commutative diagram of solid arrows
\begin{equation}
\label{diagr:valuative criterion finite type}
\begin{gathered}
\xymatrix{
\Spec(K)\ar[r] \ar[d] & X \ar[d] \\
\Spec(A) \ar[r] \ar@{.>}[ru] & \Spec(k)
}
\end{gathered}
\end{equation}
there exists at most one dotted arrow making it commutative, then $X$ is separated over $k$.
\end{lemma}

\begin{proof}
We follow the proof given in \cite[\href{https://stacks.math.columbia.edu/tag/03KV}{Tag\,03KV}]{stacks-project}. By assumption the diagonal map $\Delta:X\to X\times X$ is quasi-compact. Moreover, for any solid diagram
\begin{equation*}
\begin{gathered}
\xymatrix{
\Spec(K)\ar[r] \ar[d] & X \ar[d]^\Delta \\
\Spec(A) \ar[r] \ar@{.>}[ru] & X\times X
}
\end{gathered}
\end{equation*}
and $A$ as in the statement
there exists a dotted arrow making it commutative, namely the composition of $\Spec (A)\to X\times X$ with any of the two projections $X\times X\to X$, as these compositions coincide by assumption. Thus, by Lemma~\ref{rem:valuative criterion closedness}, $\Delta$ is universally closed, and hence a immersion as in the proof of \cite[\href{https://stacks.math.columbia.edu/tag/03KV}{Tag\,03KV}]{stacks-project}. \end{proof}

\begin{proof}[Proof of Corollary~\ref{cor:separation}]
To prove the statement, we will apply Lemma~\ref{rem:valuative criteria} above. More precisely, for every solid diagram of the form \eqref{diagr:valuative criterion finite type}
we want to show that there exists at most one dotted arrow making it commutative. Here and in the subsequent discussion, we write $M:=M_{(X,\omega),\tau}^{{\rss}}$ to simplify notation. Since $M$ is of finite type over $\C$ we may suppose moreover that in the diagram above $A$ is a discrete valuation ring over $\C$; \cf \cite[\href{https://stacks.math.columbia.edu/tag/0ARJ}{Tag\,0ARJ}]{stacks-project}.

First, we observe that by \cite[Th.\,A.8 and Rem.~A.9]{AHLH} for any morphism $f: \Spec(A)\to M$ filling in Diagram \eqref{diagr:valuative criterion finite type} there exists a lift $\tilde f: \Spec (A')\to \cX$ to the stack $\cX$, possibly after passing to a finite extension $A \hto A'$ of $A$. Next we will show that two such lifts $\tilde f,\tilde g:\Spec (A')\to \cX$ give the same morphism $\Spec(A')\to M$ after composition with the moduli morphism $\cX\to M$. Since $\Spec(A')\to\Spec(A)$ is an epimorphism in the category of algebraic spaces, this will imply the desired equality $f=g$.

Resetting notation we are led to consider a discrete valuation ring $A$ with fraction field $K$ and residue field $\C$ and two morphisms $f,g:\Spec(A)\to M$ admitting lifts $\tilde f, \tilde g$ which make the following corresponding diagrams commutative:
\begin{equation*}
\begin{gathered}
\xymatrix@=1.2cm{
\Spec(K)\ar[r] \ar[d] & \cX \ar[d] \\
\Spec(A) \ar@<0.5ex>[r]^f \ar@<-0.5ex>[r]_g \ar@<0.5ex>[ur]^(.55){\tilde f } \ar@<-0.5ex>[ru]_(.55){\tilde g}&M.
}
\end{gathered}
\end{equation*}
We remark that by Proposition \ref{prop:separation} the families of semistable sheaves that correspond to $\tilde f$ and $\tilde g$ are connected by a finite number of Langton edges, each such edge making the connection between two families of semistable sheaves on $X_{A}$. So we may reduce ourselves to the case when $\tilde f$ and $\tilde g$ correspond to the ends of a Langton edge.
By Remark \ref{rem:AHLH} we see that such an edge gives rise to a morphism $\overline{ST}_A\to\cC\!oh(X)$ that factors through the open substack $\cX\subset \cC\!oh(X)$.
As mentioned in the same remark, the projection $\overline{ST}_A\to \Spec (A)$ is a good moduli space and as such it is universal for morphisms to algebraic spaces by \cite[Th.\,6.6]{Alper13}. Consequently, the two morphisms~$f$ and~$g$ must be equal to the uniquely determined morphism $\Spec(A)\to M$ that is induced by the composition $\overline{ST}_A\to \cX\to M$.
\end{proof}

\subsection{Properness}

As a final item we have

\begin{Prop}[Properness]\label{prop:properness} The moduli space $M_{(X,\omega),\tau}^{{\rss}}$
is proper.
\end{Prop}
\begin{proof}
This follows from the analogon of Langton's valuative criterion for properness proved in \cite{TomaCriteria}.
\end{proof}

